\documentclass[aps,preprint]{revtex4}%
\usepackage{amsfonts}
\usepackage{amsmath}
\usepackage{amssymb}
\usepackage{graphicx}%
\setcounter{MaxMatrixCols}{30}
%TCIDATA{OutputFilter=latex2.dll}
%TCIDATA{Version=4.10.0.2363}
%TCIDATA{CSTFile=revtex4.cst}
%TCIDATA{Created=Sunday, January 05, 2003 12:06:13}
%TCIDATA{LastRevised=Monday, June 16, 2003 12:52:04}
%TCIDATA{<META NAME="GraphicsSave" CONTENT="32">}
%TCIDATA{<META NAME="DocumentShell" CONTENT="Articles\SW\REVTeX 4">}
%TCIDATA{Language=American English}
\newtheorem{theorem}{Theorem}
\newtheorem{acknowledgement}[theorem]{Acknowledgement}
\newtheorem{algorithm}[theorem]{Algorithm}
\newtheorem{axiom}[theorem]{Axiom}
\newtheorem{claim}[theorem]{Claim}
\newtheorem{conclusion}[theorem]{Conclusion}
\newtheorem{condition}[theorem]{Condition}
\newtheorem{conjecture}[theorem]{Conjecture}
\newtheorem{corollary}[theorem]{Corollary}
\newtheorem{criterion}[theorem]{Criterion}
\newtheorem{definition}[theorem]{Definition}
\newtheorem{example}[theorem]{Example}
\newtheorem{exercise}[theorem]{Exercise}
\newtheorem{lemma}[theorem]{Lemma}
\newtheorem{notation}[theorem]{Notation}
\newtheorem{problem}[theorem]{Problem}
\newtheorem{proposition}[theorem]{Proposition}
\newtheorem{remark}[theorem]{Remark}
\newtheorem{solution}[theorem]{Solution}
\newtheorem{summary}[theorem]{Summary}
\newenvironment{proof}[1][Proof]{\noindent\textbf{#1.} }{\ \rule{0.5em}{0.5em}}
\begin{document}
\preprint{Diagonal Operators}
\title[Diag\_Op]{A Novel Method for Generating Approximate Wavefunctions}
\author{B. Weiner}
\affiliation{Department of Physics, Pennsylvania State University, DuBois PA 15801}
\keywords{one two three}
\pacs{PACS number}

\begin{abstract}
Molecular Hamiltonians are constituted by momentum operators (e.g. kinetic
energy and interaction with electro-magnetic radiation), local potential
operators (e.g. electron-electron coulomb, electron-nucleus and
electron-external field interactions) and spin operators. In general the
effect of transformations that generate states from a reference state is to
destroy this form.

Here it is shown how one can generate all pure states by the action of a
commutative semigroup of transformations that are diagonal in the space-spin
coordinate representation that \emph{do not }destroy\emph{ }this form and thus
maintain the physical interpretability of effective Hamiltonians.

Using such transformations we formulate an approximation theory that has
physical interpretation at all orders.

\end{abstract}
\date{02/15/2003}
\startpage{1}
\maketitle
\tableofcontents


\section{Introduction\label{Intro}}

The Hamiltonian operator characterizes a physical system by describing all the
interactions effecting it and determines which states are consistent with
these interactions. Invariably physical Hamiltonians (as distinct from model
Hamiltonians) are formulated in the space-spin \emph{coordinate}
representation, which reflects the fact that the properties of natural systems
are most tangibly described primarily in terms of spatial coordinates and
geometric images in a local pointwise fashion, i.e. at any given instant the
momentum and potential energy (including the potential energy of interaction
with like particles) associated with a specific particle is expressed in terms
of a single spatial coordinate. Once the Hamiltonian has been constructed one
can form discrete matrix or even other continuous representations by the
introduction of proper bases (constituted by vectors that belong to the
Hilbert space) or generalized bases (constituted by vectors that do not belong
to the Hilbert space). However in order to obtain the physical content of such
representations one invariably relates it back to the coordinate representation.

Nearly all molecular quantum mechanical approximations to Hamiltonians and/or
states are formulated in the context of proper bases and when transformed back
to the coordinate representation produce \emph{non local }effective
Hamiltonians, which in general makes their direct physical interpretation
somewhat difficult. Consequently it would be convenient if the approximations
could be expressed directly in the coordinate representation, in which they do
have a clear physical interpretation. An added benefit of such a formalism
would be that approximations could rigorously be suggested by certain
properties of the specific system under study.

The most widely used approximation for molecular systems is the Self
Consistent Field (SCF)\ method. One of the most appealing properties of this
method is that it is based on an Independent Particle State (IPS) ansatz which
permits one to give a physical interpretation to all terms in the variational
expression and thus producing a pictorial model of processes involving atoms
and molecules. A major computational asset of this method is that it does not
involve two particle transformations thus allowing the Hamiltonian to be
expressed in terms of a fixed atomic orbital basis at all stages of the
optimization procedure i.e. the form of the Hamiltonian is preserved at all
steps. This permits one to introduce approximations to the Hamiltonian while
still maintaining the physical meaning of the various terms e.g. the coulomb
operator can be approximated by ignoring interactions between distant parts of
the molecule and core interactions replaced by pseudo potentials.

The optimal Independent Particle State (IPS)\ is found by minimizing the
expectation value of Hamiltonians with respect to (wrt) the parameters
characterizing these states. This expectation value can be alternatively
viewed as the expectation value of the exact or a fixed approximate
Hamiltonian wrt a variable IPS\ \emph{or} the expectation value of an
effective Hamiltonian that depends on the variational parameters wrt a fixed
reference IPS. However in latter case one often loses the ability to clearly
physically interpret the various terms in the effective Hamiltonian e.g. terms
that refer to local atomic centers become delocalized. When one is considering
correlated approximations beyond the SCF method this latter way of viewing the
energy expectation value is frequently more desirable, even taking into
account the loss of direct physical interpretability, as it focuses on the
most effective way of computing the energy expression e.g. consider the couple
cluster method.

The variation over trial states can be formulated in terms of transformations,
that frequently form a commuting set, acting on a fixed reference state. Even
when the set of transformations commute they invariably do not commute with
their adjoints, nor with terms in the Hamiltonian and hence produce rather
complicated effective Hamiltonians. There is however a way of combining the
desirable features of physical interpretability and fixed Hamiltonian form by
generating all pure states from a fixed independent particle reference state
by transformations that

\begin{description}
\item[$\bullet$] are diagonal in the space-spin coordinate representation

\item[$\bullet$] form self adjoint commuting sets

\item[$\bullet$] produce effective Hamiltonians that maintain the form of the
exact Hamiltonian i.e. are expressed in terms of momentum and local potentials

\item[$\bullet$] maintain the physical interpretability of all terms in
effective Hamiltonians

\item[$\bullet$] allow one to express the energy of any state in terms of the
expectation value of an effective Hamiltonian wrt a given reference state

\item[$\bullet$] produce approximations that have a direct physical interpretation.
\end{description}

In this article we describe how pure states can be generated from a reference
state by the action of abelian semigroups of operators \cite{Encyclo}, which
allows one to express the expectation value of all observable in any state in
terms of "dressed" observables in the reference state. Utilizing this
construction and considering inclusive sequences of semigroups we develop an
approximation theory. Our construction is based on the properties of
generalized bases, in particular the bases formed by the generalized
eigenvectors of the coordinate operator. In section \ref{Semi} we describe the
abelian semigroups associated with this bases and prove that one can generate
all pure states from a reference state using these semigroups. In section
\ref{EffHam} we display the effective Hamiltonian generated by the action of
this semigroup on the system Hamiltonian and in section \ref{Approx} we
develop an approximation theory based on this construction, to be followed in
section \ref{Example} with some examples and a summary in section \ref{Sum}.

\section{Generation of States\label{Semi}}

Generalized bases of Hilbert spaces have been a cornerstone of quantum
mechanics since their introduction by Dirac \cite{P.A.M.Dirac}, since their
inception they have been given a firm mathematical foundation in the context
of Distribution theory \cite{Richtmyer}, Group theory \cite{Barut} and
Topological Vector Spaces \cite{Robertson}. A full discussion of their
properties can be found in \cite{Berezin91} and they are closely related to
Rigged Hilbert spaces introduced by \cite{Bohm78}. In this section we prove
that an abelian semigroup formed by operators that are diagonal in the
space-spin generalized eigen basis generates all pure states form a cyclic
\cite{Encyclo} reference state. It is essential to consider abelian semigroups
and \emph{not} abelian groups as the latter can only maximally generate a
proper subset of states from a reference state i.e. cyclic reference states
\emph{do not} exist for these groups.

\subsection{Resolution of the Identity}

The identity operator in $\mathcal{B}\left(  \mathcal{H}^{N}\right)  $ (see
Appendix \ref{spaces} for the definition of operator spaces) can be resolved
as an integral of the projectors onto the generalized eigenvectors $\left\{
\left\vert \mathsf{z}^{N}\right\rangle \right\}  $ of the $N$-particle
space-spin position operator as%
\begin{equation}
I=\int\left\vert \mathsf{z}^{N}\right\rangle \left\langle \mathsf{z}%
^{N}\right\vert d\mathsf{z}^{N}%
\end{equation}
where $\left\vert \mathsf{z}^{N}\right\rangle =\left\vert \mathsf{z}_{1}%
\cdots\mathsf{z}_{N}\right\rangle $ is the antisymmetric tensor product of $N$
one particle generalized eigenvectors $\left\{  \left\vert \mathsf{z}%
_{j}\right\rangle |1\leq j\leq N\right\}  $ and $\mathsf{z}_{j}\equiv\left(
\mathbf{r}_{j},\xi_{j}\right)  $ for a spatial position $\mathbf{r}_{j}$ and
spin vector $\xi_{j}$. The wavefunction $\Psi\left(  \mathsf{z}^{N}\right)  $
of any pure state $\left\vert \Psi\right\rangle \left\langle \Psi\right\vert $
(see Appendix \ref{states} for a discussion of states) is given by
\begin{equation}
\Psi\left(  \mathsf{z}^{N}\right)  =\left\langle \mathsf{z}^{N}|\Psi
\right\rangle
\end{equation}
and the state vector $\left\vert \Psi\right\rangle $ has the expansion
\begin{equation}
\left\vert \Psi\right\rangle =\int\left\vert \mathsf{z}^{N}\right\rangle
\left\langle \mathsf{z}^{N}|\Psi\right\rangle d\mathsf{z}^{N}%
\end{equation}


\subsection{Hamiltonian}

The molecular spin free Hamiltonian for a system of $N$-electrons and
$M$-nuclei at fixed positions $\left\{  \mathbf{R}_{j}|1\leq j\leq M\right\}
$ is essentially local (i.e. all the terms are point interactions or are
properties of a specific point), and is defined through its tangibly physical
space-spin coordinate representation in atomic units as%
\begin{align}
H  &  =\sum_{1\leq j<k\leq M}\frac{Z_{j}Z_{k}}{\left\Vert \mathbf{R}%
_{j}-\mathbf{R}_{k}\right\Vert }\nonumber\\
&  \qquad\qquad\qquad+\int\left\{  \sum_{1\leq j\leq N}V_{ext}\left(
\mathbf{r}_{j}\right)  +\sum_{1\leq j<k\leq N}\frac{1}{\left\Vert
\mathbf{r}_{j}-\mathbf{r}_{k}\right\Vert }-\right. \nonumber\\
&  \qquad\qquad\qquad\qquad\qquad\left.  \sum_{\substack{1\leq j\leq M\\1\leq
k\leq N}}\frac{Z_{j}}{\left\Vert \mathbf{R}_{j}-\mathbf{r}_{k}\right\Vert
}-\sum_{1\leq j\leq N}\left\Vert \mathbf{p}_{j}\left(  \mathbf{r}_{j}\right)
\right\Vert ^{2}\right\}  \left\vert \mathbf{r}^{N}\right\rangle \left\langle
\mathbf{r}^{N}\right\vert d\mathbf{r}^{N} \label{hamilt1}%
\end{align}
where nuclei $j$ has atomic number $Z_{j}$ and is at the position
$\mathbf{R}_{j}$, $\mathbf{r}_{j}$ is the position of electron $j$ that has
momentum $\mathbf{p}_{j}\left(  \mathbf{r}_{j}\right)  =\nabla_{\mathbf{r}%
_{j}}-\mathbf{A}_{ext\,j}\left(  \mathbf{r}_{j}\right)  ,$ where
$\mathbf{A}_{ext\,j}\left(  \mathbf{r}_{j}\right)  $ is the sum of all
external vector potentials acting on particle $j$ and $V_{ext}\left(
\mathbf{r}_{j}\right)  $ is the sum of all external scalar potentials. For
future reference we denote the one and two particle potentials $V_{1}$ and
$V_{2}$ as
\begin{align}
V_{1}  &  =%
%TCIMACRO{\dint }%
%BeginExpansion
{\displaystyle\int}
%EndExpansion
\sum\limits_{\substack{1\leq j\leq M\\1\leq k\leq N}}\left\{  -\frac{Z_{j}%
}{\left\Vert \mathbf{R}_{j}-\mathbf{r}_{k}\right\Vert }+V_{ext}\left(
\mathbf{r}_{k}\right)  \right\}  \left\vert \mathsf{z}^{N}\right\rangle
\left\langle \mathsf{z}^{N}\right\vert d\mathsf{z}^{N}\nonumber\\
V_{2}  &  =%
%TCIMACRO{\dint }%
%BeginExpansion
{\displaystyle\int}
%EndExpansion
\left\{  \sum\limits_{1\leq j<k\leq N}\frac{1}{\left\Vert \mathbf{r}%
_{j}-\mathbf{r}_{k}\right\Vert }\right\}  \left\vert \mathsf{z}^{N}%
\right\rangle \left\langle \mathsf{z}^{N}\right\vert d\mathsf{z}^{N}%
\end{align}
and a constant internuclear interaction term
\begin{equation}
H_{0}=\sum_{1\leq j<k\leq M}\frac{Z_{j}Z_{k}}{\left\Vert \mathbf{R}%
_{j}-\mathbf{R}_{k}\right\Vert }.
\end{equation}
. By introducing a proper basis $\left\{  \left\vert \varphi_{j}\right\rangle
\right\}  $ for the one particle Hilbert space $\mathcal{H}^{1}$ one produces
a matrix representation of $H$. If this basis is constituted by localized
atomic orbitals this matrix representation still has a physically recognizable
content but if the functions $\left\{  \left\vert \varphi_{j}\right\rangle
\right\}  $ are delocalized over all centers it is hard to discern specific
physical interactions and properties. In all cases the space-spin coordinate
is the defining representation that determines the matrix one.

The energy $E\left(  \Psi\right)  $ of a state $\left\vert \Psi\right\rangle $
is given by the expectation value
\begin{equation}
E\left(  \Psi\right)  =\frac{\left\langle \Psi|H\Psi\right\rangle
}{\left\langle \Psi|\Psi\right\rangle }%
\end{equation}
where
\begin{align}
\frac{\left\langle \Psi|H\Psi\right\rangle }{\left\langle \Psi|\Psi
\right\rangle }  &  =\sum_{1\leq j<k\leq M}\frac{Z_{j}Z_{k}}{\left\Vert
\mathbf{R}_{j}-\mathbf{R}_{k}\right\Vert }+\left\langle \Psi|\Psi\right\rangle
^{-1}\left\{  \int\sum_{\substack{1\leq j\leq M\\1\leq k\leq N}}\frac
{-Z_{j}\Psi\left(  \mathsf{z}^{N}\right)  ^{\ast}\Psi\left(  \mathsf{z}%
^{N}\right)  }{\left\Vert \mathbf{R}_{j}-\mathbf{r}_{k}\right\Vert
}d\mathsf{z}^{N}\right. \nonumber\\
&  -\int\sum_{1\leq j\leq N}^{N}\left[  \left(  \nabla_{\mathbf{r}_{j}%
}-\mathbf{A}_{ext\,j}\left(  \mathsf{z}^{N}\right)  \right)  \Psi\left(
\mathsf{z}^{N}\right)  ^{\ast}\right]  \bullet\left[  \left(  \nabla
_{\mathbf{r}_{j}}-\mathbf{A}_{ext\,j}\left(  \mathsf{z}^{N}\right)  \right)
\Psi\left(  \mathsf{z}^{N}\right)  \right]  d\mathsf{z}^{N}\nonumber\\
&  +\left.  \int\sum_{1\leq j<k\leq N}\frac{\Psi\left(  \mathsf{z}^{N}\right)
^{\ast}\Psi\left(  \mathsf{z}^{N}\right)  }{\left\Vert \mathbf{r}%
_{j}-\mathbf{r}_{k}\right\Vert }d\mathsf{z}^{N}+\int V_{ext}\left(
\mathsf{z}^{N}\right)  \Psi\left(  \mathsf{z}^{N}\right)  ^{\ast}\Psi\left(
\mathsf{z}^{N}\right)  d\mathsf{z}^{N}\right\}
\end{align}
the coordinates $\mathsf{z}^{N}\equiv\mathsf{z}_{1}\cdots\mathsf{z}_{N}%
\equiv\left(  \mathbf{r}_{1},\sigma_{1}\right)  ,\cdots,\left(  \mathbf{r}%
_{N},\sigma_{N}\right)  $ and the normalization is
\begin{equation}
\left\langle \Psi|\Psi\right\rangle =\int\Psi\left(  \mathsf{z}^{N}\right)
^{\ast}\Psi\left(  \mathsf{z}^{N}\right)  d\mathsf{z}^{N}%
\end{equation}


\subsection{Pure State Generation}

In the coordinate representation, all pure states (unnormalized, see Appendix
\ref{states}) in $\mathcal{H}^{N}$ can be produced from a reference pure
state, $\left\vert \Psi_{ref}\right\rangle $, that has a wavefunction
\begin{equation}
\left\langle \mathsf{z}^{N}|\Psi_{ref}\right\rangle =\varrho\left(
\mathsf{z}^{N}\right)  e^{iS\left(  \mathsf{z}^{N}\right)  }%
\end{equation}
that is non zero $a.e.\cite{Encyclo}$ by the transformations
\begin{equation}
\left\vert \Psi\left(  \nu,\omega\right)  \right\rangle =\breve{T}\left(
\nu,\omega\right)  \left\vert \Psi_{ref}\right\rangle =\breve{\nu}%
_{N}e^{i\breve{\omega}_{N}}\left\vert \Psi_{ref}\right\rangle
\end{equation}
where the operators $\breve{T}\left(  \nu,\omega\right)  ,\breve{\nu}_{N}$ and
$\breve{\omega}_{N}$ are diagonal in the space-spin basis $\left\{  \left\vert
\mathsf{z}^{N}\right\rangle \right\}  $ and are given by
\begin{subequations}
\begin{align}
\breve{\nu}_{N}  &  =\int\nu\left(  \mathsf{z}^{N}\right)  \left\vert
\mathsf{z}^{N}\right\rangle \left\langle \mathsf{z}^{N}\right\vert
d\mathsf{z}^{N}\in\mathcal{B}\left(  \mathcal{H}^{N}\right) \\
\breve{\omega}_{N}  &  =\int\omega\left(  \mathsf{z}^{N}\right)  \left\vert
\mathsf{z}^{N}\right\rangle \left\langle \mathsf{z}^{N}\right\vert
d\mathsf{z}^{N}\in\mathcal{B}\left(  \mathcal{H}^{N}\right)
\end{align}
and form a commuting set, which we denote as $\mathcal{C}_{c}$. $\lceil$A
notational aside: In this article when it is important to distinguish between
an operator belonging to $\mathcal{B}\left(  \mathcal{H}^{N}\right)  $ and its
integral kernel, we denote the operator by $\breve{T}$ and the integral kernel
by $T\left(  \mathsf{z}^{N}\right)  $ or $T\left(  \mathsf{z}^{N}%
,\mathsf{z}^{\prime N}\right)  ,$ if it not important to distinguish between
we just use $T\rfloor$.

The strong closure \cite{Richtmyer} of the space $\mathcal{C}_{c}$ of diagonal
integral operators that map $\mathcal{H}^{N}\rightarrow\mathcal{H}^{N}$ is the
space of integral operators with symmetric integral kernels $\left\{  T\left(
\mathsf{z}^{N}\right)  \right\}  $ that belong to the Banach space $L_{\infty
}\left(  \mathsf{Z}^{N}\right)  $ $\cite{Encyclo}$, and form a semigroup which
is in particular a $w^{\ast}$-algebra. \cite{Thirring3}. The following extract
from Thirring \cite{Thirring3}\textit{ }should be noted:
\end{subequations}
\begin{quotation}
$L_{\infty}\left(  \mathsf{Z}^{N}\right)  ,$ \textit{considered as
multiplication operators on} $L_{2}\left(  \mathsf{Z}^{N}\right)  $ \textit{is
maximally abelian. }\emph{Every function in }$L_{2}\left(  \mathsf{Z}%
^{N}\right)  $\emph{\ that is non zero }$a.e.$\emph{\ is a cyclic
vector}\textit{. Functions that vanish on some set} $\left[  \Omega
^{N}\right]  ^{C}=\mathsf{Z}^{N}-\Omega^{N}\subset\mathsf{Z}^{N}$ \textit{form
invariant subspaces. Thus} $L_{\infty}\left(  \mathsf{Z}^{N}\right)  $
\textit{is reducible and not a factor }\cite{Encyclo}.
\end{quotation}

One can see that $T\left(  \nu,\omega\right)  \left(  \mathsf{z}^{N}\right)  $
(the integral kernel of $\breve{T}\left(  \nu,\omega\right)  )$ is a
parameterization of a representation of an abelian $w^{\ast}$-algebra,
$\mathfrak{A}$, on $L_{2}\left(  \mathsf{Z}^{N}\right)  $ by $L_{\infty
}\left(  \mathsf{Z}^{N}\right)  $. This representation is reducible as
$L_{2}\left(  \mathsf{Z}^{N}\right)  $ contains the following invariant
subspaces
\begin{equation}
L_{2}\left(  \Omega^{N}\right)  =\left\{  \Psi\left\vert \Psi\left(
\mathsf{z}^{N}\right)  =0\,\text{if \thinspace}\mathsf{z}^{N}\notin\Omega
^{N}\,\text{,\thinspace}\int_{\Omega^{N}}d\mathsf{z}^{N}\neq
0\,\text{and\thinspace}\Omega^{N}\subset\mathsf{Z}^{N}\right.  \right\}
\end{equation}
The operator $\breve{T}\left(  \nu,\omega\right)  $ produces elements that
have amplitudes that belong to $L_{2}\left(  \Omega^{N}\right)  $ when
$\operatorname{Supp}\left(  \nu\right)  =$ $\Omega^{N}$ $\cite{Richtmyer}$.

The set of invertible diagonal operators in $\mathfrak{A}$ form a group,
$\mathfrak{G}$, that also has a reducible representation on $L_{\infty}\left(
\mathsf{Z}^{N}\right)  $, as clearly $L_{2}\left(  \Omega^{N}\right)  $ are
also invariant subspaces of $\mathfrak{G}$. However \emph{it is not true} that
every function in $L_{2}\left(  \mathsf{Z}^{N}\right)  $\ that is nonzero
$a.e.$\ is a cyclic vector for $\mathfrak{G}$ and in fact there do not exist
\emph{any} cyclic vectors for $\mathfrak{G}$ in $L_{2}\left(  \mathsf{Z}%
^{N}\right)  $ and thus one can never generate all pure states by the action
of $\mathfrak{G}$ on a reference state.

We can contrast this semigroup to ones formed by proper bases of $N$-particle
space by comparing the preceding to the discrete case by focusing on a proper
basis $\left\{  \left\vert \Phi_{j}\right\rangle \right\}  $ of $\mathcal{H}%
^{N}$, the abelian algebra, $\mathcal{C}_{d},$ generated by $\left\{
%TCIMACRO{\tsum \limits_{1\leq j\leq\infty}}%
%BeginExpansion
{\textstyle\sum\limits_{1\leq j\leq\infty}}
%EndExpansion
\lambda_{j}\left\vert \Phi_{j}\right\rangle \left\langle \Phi_{j}\right\vert
\right\}  $ and a reference state $\left\vert \Phi_{ref}\right\rangle =%
%TCIMACRO{\tsum \limits_{1\leq j\leq\infty}}%
%BeginExpansion
{\textstyle\sum\limits_{1\leq j\leq\infty}}
%EndExpansion
\alpha_{j}\left\vert \Phi_{j}\right\rangle $ that has the property that
$\alpha_{j}\neq0\,\forall i$ $\left(  \text{see Appendix \ref{Difference}%
}\right)  $. Similarly to $\mathcal{C}_{c}$ the abelian algebra $\mathcal{C}%
_{d}$ can also generate all pure states from such a reference state but it
\emph{does not }have such desirable properties wrt $H$ as $\mathcal{C}_{c}$
does e.g. the potential terms in $H$ are no longer diagonal nor will
$\mathcal{C}_{c}$ commute with these terms.

\section{Effective Hamiltonian\label{EffHam}}

Any pure state $N$-particle density operator $\breve{D}^{N}\left(  \nu
,\omega\right)  $, (in general unnormalized), can be generated from a cyclic
reference density $\mathcal{\breve{D}}_{ref}^{N}$ by the transformation
\begin{equation}
\mathcal{\breve{D}}^{N}\left(  \nu,\omega\right)  =\hat{T}\left(
\mathcal{\breve{D}}_{ref}^{N}\right)  =\breve{T}\left(  \nu,\omega\right)
\mathcal{\breve{D}}_{ref}^{N}\breve{T}\left(  \nu,\omega\right)  ^{\dagger}%
\end{equation}
The transformation $\hat{T}\left(  \nu,\omega\right)  $ is positive as it maps
the set of positive operators $\mathcal{B}_{1+}\left(  \mathcal{H}^{N}\right)
$ to itself and the set of all such $\hat{T}$ correspond to the semigroup
$\mathcal{C}_{c}$.

As the energy expectation value $E\left(  \nu,\omega\right)  $ can be
expressed as
\begin{equation}
E\left(  \nu,\omega\right)  =\frac{Tr\left\{  H\mathcal{\breve{D}}^{N}\left(
\nu,\omega\right)  \right\}  }{Tr\left\{  \mathcal{\breve{D}}^{N}\left(
\nu,\omega\right)  \right\}  }=\frac{Tr\left\{  H\breve{T}\left(  \nu
,\omega\right)  \mathcal{\breve{D}}_{ref}^{N}\breve{T}\left(  \nu
,\omega\right)  ^{\dagger}\right\}  }{Tr\left\{  \breve{T}\left(  \nu
,\omega\right)  \mathcal{\breve{D}}_{ref}^{N}\breve{T}\left(  \nu
,\omega\right)  ^{\dagger}\right\}  }=\frac{Tr\left\{  \breve{T}\left(
\nu,\omega\right)  ^{\dagger}H\breve{T}\left(  \nu,\omega\right)
\mathcal{\breve{D}}_{ref}^{N}\right\}  }{Tr\left\{  \breve{T}\left(
\nu,\omega\right)  \mathcal{\breve{D}}_{ref}^{N}\breve{T}\left(  \nu
,\omega\right)  ^{\dagger}\right\}  }%
\end{equation}
the action of the transformation $\hat{T}\left(  \nu,\omega\right)  $ can be
incorporated into an effective Hamiltonian $H\left(  \nu,\omega\right)  $
defined by
\begin{equation}
H\left(  \nu,\omega\right)  =\breve{T}\left(  \nu,\omega\right)  ^{\dagger
}H\breve{T}\left(  \nu,\omega\right)
\end{equation}
and all expectation values can be taken wrt the IPS $\mathcal{\breve{D}}%
_{ref}^{N}$.

The Hamiltonians, $H$, we consider are of the form%
\begin{equation}
H=H_{0}-\mathbf{P\bullet P}+V_{1}+V_{2}%
\end{equation}
where $H_{0}$ is a constant, $\mathbf{P}$ the total \emph{generalized
}momentum operator for the system, $V_{1}$ is the sum of all one particle
potential operators and $V_{2}$ is the sum of all two particle potential
operators. The \emph{generalized }momentum is defined in terms of one particle
\emph{generalized }momentum operators as
\begin{equation}
\mathbf{P=}%
%TCIMACRO{\dsum \limits_{1\leq j\leq N}}%
%BeginExpansion
{\displaystyle\sum\limits_{1\leq j\leq N}}
%EndExpansion
\mathbf{p}_{j}%
\end{equation}
where
\begin{equation}
\mathbf{p}_{j}=-i\mathbf{\nabla}_{j}-\mathbf{A}_{j}%
\end{equation}
$\mathbf{A}_{j}$ is the total external vector potential acting on particle $j$
and
\begin{equation}
-i\mathbf{\nabla}_{j}=-i\left(
\begin{array}
[c]{c}%
\frac{\partial}{\partial x_{j}}\\
\frac{\partial}{\partial y_{j}}\\
\frac{\partial}{\partial z_{j}}%
\end{array}
\right)  ;\quad\mathbf{A}_{j}=\left(
\begin{array}
[c]{c}%
A_{jx}\\
A_{jy}\\
A_{jz}%
\end{array}
\right)
\end{equation}
The potentials $\left\{  \mathbf{A}_{j}\right\}  $, $V_{1}$ and $V_{2}$ are
all diagonal in the space-spin representation. The operators $\left\{
\breve{T}\left(  \nu,\omega\right)  \right\}  $ \emph{all} \emph{commute }with
these potentials, which is of course by design, so that
\begin{align}
\breve{T}\left(  \nu,\omega\right)  ^{\dagger}V\breve{T}\left(  \nu
,\omega\right)   &  =\breve{\nu}_{N}e^{-i\breve{\omega}_{N}}V\breve{\nu}%
_{N}e^{i\breve{\omega}_{N}}=\left(  \breve{\nu}_{N}\right)  ^{2}V\nonumber\\
\breve{T}\left(  \nu,\omega\right)  ^{\dagger}\mathbf{A}_{j}\breve{T}\left(
\nu,\omega\right)   &  =\breve{\nu}_{N}e^{-i\breve{\omega}_{N}}\mathbf{A}%
_{j}\breve{\nu}_{N}e^{i\breve{\omega}_{N}}=\left(  \breve{\nu}_{N}\right)
^{2}\mathbf{A}_{j}%
\end{align}
where $V=V_{1}$ $+$ $V_{2}$
\begin{align}
\left(  \breve{\nu}_{N}\right)  ^{2}V  &  =\int\nu\left(  \mathsf{z}%
^{N}\right)  ^{2}V\left(  \mathsf{z}^{N}\right)  \left\vert \mathsf{z}%
^{N}\right\rangle \left\langle \mathsf{z}^{N}\right\vert d\mathsf{z}%
^{N}\nonumber\\
\left(  \breve{\nu}_{N}\right)  ^{2}A_{j\kappa}  &  =\int\nu\left(
\mathsf{z}^{N}\right)  ^{2}A_{j\kappa}\left(  \mathsf{z}^{N}\right)
\left\vert \mathsf{z}^{N}\right\rangle \left\langle \mathsf{z}^{N}\right\vert
d\mathsf{z}^{N};\qquad\kappa=x,y,z
\end{align}
i.e. the potentials are scaled by the integral kernels $\left\{  \nu\left(
\mathsf{z}^{N}\right)  \right\}  $. It is \emph{important }to note that while
the Hamiltonian $H$ only contains potentials describing up to two particle
interactions, the integral kernels $\left\{  \nu\left(  \mathsf{z}^{N}\right)
\right\}  $ are in general \emph{multiparticle} functions and thus produce
multiparticle potentials $\left\{  \nu\left(  \mathsf{z}^{N}\right)
^{2}V\left(  \mathsf{z}^{N}\right)  \right\}  $ describing up to simultaneous
$N$-particle interactions. Furthermore the $N$-particle local scaling factors
$\nu\left(  \mathsf{z}^{N}\right)  $ contain phase information pertaining to
$1,2,\cdots,M<N$ particle interactions and the $N$-particle local phase
$\omega\left(  \mathsf{z}^{N}\right)  $ contains local scaling information
pertaining to $1,2,\cdots,M<N$ particle interactions as will discussed in the examples.

The effect of the unitary part $\left\{  e^{i\omega_{N}}\right\}  $ of the
transformations $\left\{  \breve{T}\left(  \nu,\omega\right)  \right\}  $ on
the momentum operators $\left\{  -i\nabla_{j}\right\}  $ (for simplicity we
consider the case where there is no external vector potential) is to produce a
generalized gauge transformation i.e. producing the component operators
\begin{equation}
e^{-i\breve{\omega}_{N}}\left(  -i\frac{\partial}{\partial x_{j}}\right)
e^{i\breve{\omega}_{N}}=-i\frac{\partial}{\partial x_{j}}+\frac{\partial
\breve{\omega}_{N}}{\partial x_{j}}%
\end{equation}
and $1$-particle vector operators
\begin{equation}
e^{-i\breve{\omega}_{N}}\left(  -i\nabla_{j}\right)  e^{i\breve{\omega}_{N}%
}=-i\nabla_{j}+\nabla_{j}\breve{\omega}_{N}%
\end{equation}
(again note that $\breve{\omega}_{N}$ are $N$-particle operators). The effect
of the scaling $\breve{\nu}_{N}$ on this is
\begin{equation}
\left(  -i\nabla_{j}+\nabla_{j}\breve{\omega}_{N}\right)  \breve{\nu}%
_{N}=-i\breve{\nu}_{N}\nabla_{j}-i\left(  \nabla_{j}\breve{\nu}_{N}\right)
+\breve{\nu}_{N}\left(  \nabla_{j}\breve{\omega}_{N}\right)
\end{equation}
leading to an effective kinetic energy operator (see Appendix \ref{KinEnTrans}%
) which is a sum of kinetic, momentum and potential terms
\begin{align}
&  K\left(  \nu,\omega\right)  =\breve{T}\left(  \nu,\omega\right)  ^{\dagger
}\left(  \mathbf{P\bullet P}\right)  \breve{T}\left(  \nu,\omega\right)
=\nonumber\\
&  \qquad-\nabla\bullet\left(  \breve{\nu}^{2}\right)  \mathbb{I}\bullet
\nabla-\nabla\bullet\left(  \breve{\nu}_{N}^{2}\nabla\breve{\omega}%
_{N}-i\breve{\nu}_{N}\nabla\breve{\nu}_{N}\right) \nonumber\\
&  \qquad\qquad-\left(  \breve{\nu}_{N}^{2}\nabla\breve{\omega}_{N}%
-i\breve{\nu}_{N}\nabla\breve{\nu}_{N}\right)  \bullet\nabla+\left(
\breve{\nu}_{N}\nabla\breve{\omega}_{N}-i\nabla\breve{\nu}_{N}\right)
\bullet\mathbb{I}\bullet\left(  \breve{\nu}_{N}\nabla\breve{\omega}%
_{N}-i\nabla\breve{\nu}_{N}\right)  \label{effkin}%
\end{align}
where $\mathbb{I}$ is $3N\times3N$ unit matrix. Thus the effective Hamiltonian
is given by%
\begin{equation}
H\left(  \nu,\omega\right)  =H_{0}+K\left(  \nu,\omega\right)  +\left(
\breve{\nu}_{N}\right)  ^{2}V \label{effham2}%
\end{equation}
which displays that the effective scalar potential does not depend on the
$N$-particle local phase factor $\exp\left\{  i\omega\left(  \mathsf{z}%
^{N}\right)  \right\}  $, but it should be noted that it \emph{does }depend on
phase relationships between the individual particles.

\section{Approximations\label{Approx}}

One can expand the transformation $\breve{T}\left(  \nu,\omega\right)
\in\mathcal{B}\left(  \mathcal{H}^{N}\right)  $ in terms of antisymmetrized
tensor products of transformations \cite{Marcus}\cite{Greub} $\breve{T}%
_{n_{j}}\in\mathcal{B}\left(  \mathcal{H}^{n_{j}}\right)  $ associated with
partitions $\left\{  n_{1},\ldots,n_{m}\right\}  $ of $N$, i.e. for sets of
$\left\{  \left.  n_{j}\right\vert 1\leq j\leq m\right\}  $ where $0\leq
n_{j}\leq N$ and $%
%TCIMACRO{\tsum \limits_{1\leq j\leq m}}%
%BeginExpansion
{\textstyle\sum\limits_{1\leq j\leq m}}
%EndExpansion
n_{j}=N$, as
\begin{equation}
\breve{T}\left(  \nu,\omega\right)  =%
%TCIMACRO{\dsum \limits_{1\leq m\leq N}}%
%BeginExpansion
{\displaystyle\sum\limits_{1\leq m\leq N}}
%EndExpansion
\,%
%TCIMACRO{\dsum \limits_{\left\{  n_{1},\ldots,n_{m}\right\}  }}%
%BeginExpansion
{\displaystyle\sum\limits_{\left\{  n_{1},\ldots,n_{m}\right\}  }}
%EndExpansion
c_{n_{1}\ldots n_{m}}\breve{T}_{n_{1}}\left(  \nu_{n_{1}},\omega_{n_{1}%
}\right)  \wedge\cdots\wedge\breve{T}_{n_{m}}\left(  \nu_{n_{m}},\omega
_{n_{m}}\right)  \label{Transexp}%
\end{equation}
where $\nu_{n_{p}}$ defines the operator $\breve{\nu}_{n_{p}}$ through an
integral kernel $\nu_{n_{p}}\left(  \mathsf{z}^{n_{p}}\right)  $ with
variables $\mathsf{z}^{n_{p}}\equiv\mathsf{z}_{q_{p}+1}\cdots\mathsf{z}%
_{q_{p}+n_{p}}$, where $q_{p}=n_{1}+\cdots+n_{p-1}$. The antisymmetrized
product operator $\breve{T}_{n_{1}}\left(  \nu_{n_{1}},\omega_{n_{1}}\right)
\wedge\cdots\wedge\breve{T}_{n_{m}}\left(  \nu_{n_{m}},\omega_{n_{m}}\right)
$ is defined by
\begin{align}
&  \breve{T}_{n_{1}}\left(  \nu_{n_{1}},\omega_{n_{1}}\right)  \wedge
\cdots\wedge\breve{T}_{n_{m}}\left(  \nu_{n_{m}},\omega_{n_{m}}\right)
\left\vert \Phi_{J_{n_{1}}}^{n_{1}}\wedge\cdots\wedge\Phi_{J_{n_{m}}}^{n_{m}%
}\right\rangle \nonumber\\
&  \qquad\qquad\qquad\qquad\qquad\qquad\qquad\qquad=\left\vert \breve
{T}_{n_{1}}\left(  \nu_{n_{1}},\omega_{n_{1}}\right)  \Phi_{J_{n_{1}}}^{n_{1}%
}\wedge\cdots\wedge\breve{T}_{n_{m}}\left(  \nu_{n_{m}},\omega_{n_{m}}\right)
\Phi_{J_{n_{m}}}^{n_{m}}\right\rangle
\end{align}
for $1\leq J_{n_{j}}\leq\binom{r}{n_{j}},\,1\leq j\leq m$ for a specific fixed
partition $\left\{  n_{1},\ldots,n_{m}\right\}  $ where $\left\{  \left.
\left\vert \Phi_{J_{n_{j}}}^{n_{j}}\right\rangle \right\vert 1\leq J_{n_{j}%
}\leq\binom{r}{n_{j}};\left\{  n_{1},\ldots,n_{m}\right\}  \right\}  $ are
proper bases for $\left\{  \left.  \mathcal{H}^{n_{j}}\right\vert 1\leq j\leq
m\right\}  $. If a factor $\breve{T}_{n_{p}}\left(  \nu_{n_{p}},\omega_{n_{p}%
}\right)  $ can be expressed as a product of lower order factors e.g.
\begin{equation}
\breve{T}_{n_{p}}\left(  \nu_{n_{p}},\omega_{n_{p}}\right)  =\breve{T}_{p_{1}%
}\left(  \nu_{p_{1}},\omega_{p_{1}}\right)  \wedge\cdots\wedge\breve{T}%
_{p_{m}}\left(  \nu_{p_{m}},\omega_{p_{m}}\right)
\end{equation}
where $%
%TCIMACRO{\tsum \limits_{1\leq j\leq m}}%
%BeginExpansion
{\textstyle\sum\limits_{1\leq j\leq m}}
%EndExpansion
p_{j}=n_{p}$ for a partition $\left\{  p_{1},\ldots,p_{m}\right\}  $ we say it
is decomposable if not, i.e. it can only be expressed as a sum of such
products, we say it is indecomposable.

In general it is possible to express the expansion eq. $\left(  \ref{Transexp}%
\right)  $ of $\breve{T}\left(  \nu,\omega\right)  $ in terms of
\emph{antisymmetric} tensor products $\breve{T}_{n_{1}}\left(  \nu_{n_{1}%
},\omega_{n_{1}}\right)  \wedge\cdots\wedge\breve{T}_{n_{m}}\left(  \nu
_{n_{m}},\omega_{n_{m}}\right)  $ of indecomposable operators $\left\{
\breve{T}_{n_{p}}\left(  \nu_{n_{p}},\omega_{n_{p}}\right)  \right\}  $ whose
integral kernels are symmetric functions $\left(  \text{see Appendix
}\ref{symker}\right)  $ given by the \emph{symmetric }\cite{Marcus}%
\cite{Greub} tensor product of the kernels of these operators as\emph{ }
\begin{equation}
T_{n_{1}\cdots n_{m}}\left(  \mathsf{z}^{N}\right)  =T_{n_{1}}\left(
\mathsf{z}^{n_{1}}\right)  \vee\cdots\vee T_{n_{m}}\left(  \mathsf{z}^{n_{m}%
}\right)
\end{equation}


Approximations to the set of transformations $\left\{  \breve{T}\left(
\nu,\omega\right)  \right\}  $ can be generated by placing restrictions on the
expansion eq. $\left(  \ref{Transexp}\right)  $ to various classes of
partitions $\left\{  n_{1},\ldots,n_{m}\right\}  $, the types of operators in
the products and by limiting the choice of integral kernels $\nu\left(
\mathsf{z}^{N}\right)  ,\omega\left(  \mathsf{z}^{N}\right)  $ by expanding
them in terms of fixed sets of function. For example some different types of
products are produced by:

\begin{enumerate}
\item The partition $\left\{  1,\ldots,1\right\}  $ that leads to one particle
operators acting in $N$-particle space $\mathcal{H}^{N}$ whose integral
kernels are formed by a product of identical one particle integral kernels
\begin{equation}
T\left(  \mathsf{z}^{N}\right)  =T_{1}^{N}\left(  \mathsf{z}^{N}\right)
=T_{1}\left(  \mathsf{z}_{1}\right)  \cdots T_{1}\left(  \mathsf{z}%
_{N}\right)
\end{equation}
and corresponds to a linear transformations of $N$-particle space
$\mathcal{H}^{N}$ through%
\begin{equation}
T_{1}\wedge\cdots\wedge T_{1}\left\vert \varphi_{j_{1}}\cdots\varphi_{j_{N}%
}\right\rangle =\left\vert T_{1}\varphi_{j_{1}}\cdots T_{1}\varphi_{j_{N}%
}\right\rangle ;\quad1\leq j_{1}<\cdots<j_{N}\leq r
\end{equation}
where $\left\{  \left.  \left\vert \varphi_{j}\right\rangle \right\vert 1\leq
j\leq r\right\}  $ forms a complete orthonormal basis for $\mathcal{H}^{1}$ of
dimension $r$.

\item The partition $\left\{  1,\ldots,1\right\}  $ that leads to $N$-particle
operators acting in $N$-particle space $\mathcal{H}^{N}$ whose integral
kernels are formed by a product of different one particle integral kernels
\begin{equation}
T\left(  \mathsf{z}^{N}\right)  =\sum_{\sigma\in\mathcal{S}_{N}}T_{1}\left(
\mathsf{z}_{\sigma\left(  1\right)  }\right)  \cdots T_{N}\left(
\mathsf{z}_{\sigma\left(  N\right)  }\right)
\end{equation}
These transformations correspond to
\begin{equation}
T_{1}\wedge\cdots\wedge T_{N}\left\vert \varphi_{1}\cdots\varphi
_{N}\right\rangle =\left\vert T_{1}\varphi_{1}\cdots T_{N}\varphi
_{N}\right\rangle \qquad\forall\left\vert \varphi_{j}\right\rangle
\in\mathcal{H}^{N},1\leq j\leq N
\end{equation}
which\emph{\ are not induced }by linear transformations of one particle space
$\mathcal{H}^{1}$ unless $T_{1}=\cdots=T_{N}$. The integral kernel of this
transformation can be developed using the polar decomposition $T_{j}\left(
\mathsf{z}\right)  =\nu_{j}\left(  \mathsf{z}\right)  e^{i\omega_{j}\left(
\mathsf{z}\right)  }$ to give
\begin{align}
T\left(  \mathsf{z}^{N}\right)   &  =\sum_{\sigma\in\mathcal{S}_{N}}%
\nu_{\sigma\left(  1\right)  }\left(  \mathsf{z}_{1}\right)  e^{i\omega
_{\sigma\left(  1\right)  }\left(  \mathsf{z}_{1}\right)  }\cdots\nu
_{\sigma\left(  N\right)  }\left(  \mathsf{z}_{N}\right)  e^{i\omega
_{\sigma\left(  N\right)  }\left(  \mathsf{z}_{N}\right)  }\nonumber\\
&  =\sum_{\sigma\in\mathcal{S}_{N}}\nu_{\sigma\left(  1\right)  }\left(
\mathsf{z}_{1}\right)  \cdots\nu_{\sigma\left(  N\right)  }\left(
\mathsf{z}_{N}\right)  e^{i\left\{  \omega_{\sigma\left(  1\right)  }\left(
\mathsf{z}_{1}\right)  +\cdots+\omega_{\sigma\left(  N\right)  }\left(
\mathsf{z}_{N}\right)  \right\}  }%
\end{align}
One should be careful to note that in this expansion $\nu\left(
\mathsf{z}^{N}\right)  \neq\sum_{\sigma\in\mathcal{S}_{N}}\nu_{\sigma\left(
1\right)  }\left(  \mathsf{z}_{1}\right)  \cdots\nu_{\sigma\left(  N\right)
}\left(  \mathsf{z}_{N}\right)  ,$ and $\omega\left(  \mathsf{z}^{N}\right)
\neq\sum\limits_{1\leq j\leq N}\omega_{\sigma\left(  j\right)  }\left(
\mathsf{z}_{j}\right)  $ but each are in fact complicated functions of
$\left\{  \nu_{j},\omega_{j}|1\leq j\leq N\right\}  $. Thus although $T\left(
\mathsf{z}^{N}\right)  =\nu\left(  \mathsf{z}^{N}\right)  e^{i\omega\left(
\mathsf{z}^{N}\right)  }$ and the norm $\left\Vert \left\vert T_{1}\varphi
_{1}\cdots T_{N}\varphi_{N}\right\rangle \right\Vert $ is \emph{independent}
of\emph{ }the\emph{ }N-particle phase $\omega\left(  \mathsf{z}^{N}\right)  $
it is \emph{not }independent of the individual particle phases $\left\{
\omega_{j}\left(  \mathsf{z}_{j}\right)  \right\}  $. Another important point
to note from this example is that IPSs \emph{are not all} generated from a
reference IPS\ by one particle diagonal operators i.e. by transformations of
the form $T_{1}\wedge\cdots\wedge T_{1}$ but they \emph{are all }generated by
$N$-particle transformations of the form $T_{1}\wedge T_{2}\wedge\cdots\wedge
T_{N},$ (see section \ref{NVCT} for a discussion of how these diagonal
$N$-particle operators correspond to non diagonal $1$-particle operators)

\item The partition $\left\{  2,\ldots,2\right\}  $ that leads to $N$-particle
operators acting in $N$-particle space $\mathcal{H}^{N}$ whose integral
kernels are formed by a product of identical two particle integral kernels (of
course there must be an even number of particles in this case)
\begin{equation}
T\left(  \mathsf{z}^{N}\right)  =T_{2}^{\frac{N}{2}}\left(  \mathsf{z}%
^{N}\right)  =\sum_{\sigma\in\mathcal{S}_{N}}T_{2}\left(  \mathsf{z}%
_{\sigma\left(  1\right)  }\mathsf{z}_{\sigma\left(  2\right)  }\right)
\cdots T_{2}\left(  \mathsf{z}_{\sigma\left(  N-1\right)  }\mathsf{z}%
_{\sigma\left(  N\right)  }\right)
\end{equation}
through the definition%
\begin{equation}
T_{2}\wedge\cdots\wedge T_{2}\left\vert g_{1}\cdots g_{\frac{N}{2}%
}\right\rangle =\left\vert T_{2}g_{1}\cdots T_{2}g_{\frac{N}{2}}\right\rangle
\qquad\forall\left\vert g_{j}\right\rangle \in\mathcal{H}^{2},1\leq j\leq
\frac{N}{2}%
\end{equation}

\end{enumerate}

Using the expansion eq. $\left(  \ref{Transexp}\right)  $ one can decompose
$T\left(  \mathsf{z}^{N}\right)  $ in terms of indecomposable parts indexed by
$\left\{  n_{1},\ldots,n_{m}\right\}  $ e.g.
\begin{equation}
T\left(  \mathsf{z}^{N}\right)  =\sum_{1\leq m\leq N}\,%
%TCIMACRO{\tsum \limits_{\left\{  n_{1},\ldots,n_{m}\right\}  }}%
%BeginExpansion
{\textstyle\sum\limits_{\left\{  n_{1},\ldots,n_{m}\right\}  }}
%EndExpansion
T_{n_{1}\ldots n_{m}}\left(  \mathsf{z}^{N}\right)  \label{vwexpan}%
\end{equation}
and formulate approximations to the energy functional $E\left(  \nu
,\omega\right)  $ by truncating and approximating terms in this expansion. For
instance one might choose the parts indexed by $\left\{  1,\cdots,1\right\}  $
which contains integral kernels of the form
\begin{equation}
T\left(  \mathsf{z}^{N}\right)  =T_{1}\left(  \mathsf{z}_{1}\right)
\vee\cdots\vee T_{N}\left(  \mathsf{z}_{N}\right)
\end{equation}
which includes $p$-particle operator integral kernels of the form%
\begin{equation}
T_{1}\left(  \mathsf{z}_{1}\right)  \vee\cdots\vee T_{p}\left(  \mathsf{z}%
_{p}\right)  \vee\delta\left(  \mathsf{z}_{p+1}\right)  \vee\cdots\vee
\delta\left(  \mathsf{z}_{N}\right)  ;\quad T_{j}\neq T_{k},1\leq j<k\leq
p;1\leq p\leq N
\end{equation}
and $1$-particle operators integral kernels of the form
\begin{equation}
T_{1}\left(  \mathsf{z}_{1}\right)  \vee\cdots\vee T_{1}\left(  \mathsf{z}%
_{N}\right)  . \label{symprod}%
\end{equation}
or just restrict attention to the operator kernels of eq. $\left(
\ref{symprod}\right)  $. A higher level of approximation would include parts
indexed by $\left\{  1,\cdots,1\right\}  $ and all or a selection from
integral kernels indexed by $\left\{  2,1,\cdots,1\right\}  ,\left\{
2,2,1,\cdots,1\right\}  ,\ldots,\left\{  2,2,\cdots,2\right\}  $ and so on....

\section{Commutative versus Noncommutative Transformations \label{NVCT}}

Example 2 of section \ref{Approx} showed that all IPSs can be generated by an
\emph{abelian} semigroup $\mathcal{C}_{c}$ of diagonal $N$-particle operators,
$T_{1}\wedge\cdots\wedge T_{N}$, composed of $N$ diagonal one particle
operators $\left\{  \left.  T_{j}\right\vert 1\leq j\leq N\right\}  $. At the
same time it is well known that all IPSs can be generated by the \emph{non
abelian} general linear group, $\mathcal{GL}\left(  \mathcal{H}^{1}\right)  $,
of transformations that map $\mathcal{H}^{1}\rightarrow\mathcal{H}^{1}$. In
this section we explicitly describe a correspondence between $\mathcal{C}_{c}$
and $\mathcal{GL}\left(  \mathcal{H}^{1}\right)  $ by considering the
equivalence of $N$-dimensional subspaces, $V\left(  \mathbf{\varphi}\right)
,$ spanned by the vectors $\left\{  \left.  \left\vert \varphi_{j}%
\right\rangle \right\vert 1\leq j\leq N\right\}  $ of $\mathcal{H}^{1}$ to
$N$-particle IPSs $\left\vert \varphi_{1}\cdots\varphi_{N}\right\rangle $. The
diagonal transformation
\begin{equation}
T_{1}\wedge\cdots\wedge T_{N}\left\vert \varphi_{1}\cdots\varphi
_{N}\right\rangle =\left\vert T_{1}\varphi_{1}\cdots T_{N}\varphi
_{N}\right\rangle
\end{equation}
can be associated with the transformation $\mathcal{T}\left(  \mathbf{\varphi
}\right)  :V\left(  \mathbf{\varphi}\right)  \rightarrow V_{T}\left(
\mathbf{\varphi}\right)  $ between the subspaces $V\left(  \mathbf{\varphi
}\right)  $ spanned by the \emph{linearly independent }vectors $\left\{
\left.  \left\vert \varphi_{j}\right\rangle \right\vert 1\leq j\leq N\right\}
$ and $V_{T}\left(  \mathbf{\varphi}\right)  $ spanned by the
\emph{orthonormal} vectors $\left\{  \left.  \left\vert \psi_{j}\right\rangle
\right\vert 1\leq j\leq N\right\}  $ defined by%
\begin{equation}
\mathcal{T}\left(  \mathbf{\varphi}\right)  \left(  \left\vert \varphi
_{1}\right\rangle ,\cdots,\left\vert \varphi_{N}\right\rangle \right)
=\left(  \left\vert \psi_{1}\right\rangle ,\cdots,\left\vert \psi
_{N}\right\rangle \right)  =\left(  \left\vert T_{1}\varphi_{1}\right\rangle
,\cdots,\left\vert T_{N}\varphi_{N}\right\rangle \right)  \mathbf{\Delta
}\left(  \mathbf{\varphi}\right)  ^{-\frac{1}{2}}%
\end{equation}
where $\Delta\left(  \mathbf{\varphi}\right)  _{jk}=\left\langle T_{j}%
\varphi_{j}|T_{k}\varphi_{k}\right\rangle $ and
\begin{equation}
\Delta\left(  \mathbf{\varphi}\right)  \left(  \left\vert T_{1}\varphi
_{1}\right\rangle ,\cdots,\left\vert T_{N}\varphi_{N}\right\rangle \right)
=\left(  \left\vert T_{1}\varphi_{1}\right\rangle ,\cdots,\left\vert
T_{N}\varphi_{N}\right\rangle \right)  \mathbf{\Delta}\left(  \mathbf{\varphi
}\right)  \label{delta}%
\end{equation}
i.e.%
\begin{equation}
\left\vert \psi_{j}\right\rangle =%
%TCIMACRO{\tsum \limits_{1\leq k\leq N}}%
%BeginExpansion
{\textstyle\sum\limits_{1\leq k\leq N}}
%EndExpansion
\left\vert T_{k}\varphi_{k}\right\rangle \left(  \mathbf{\Delta}\left(
\mathbf{\varphi}\right)  ^{-\frac{1}{2}}\right)  _{kj};\qquad1\leq j\leq N
\end{equation}
Using eq. $\left(  \ref{delta}\right)  $ we obtain
\begin{equation}
\left\vert \psi_{1}\cdots\psi_{N}\right\rangle =\det\left(  \mathbf{\Delta
}\left(  \mathbf{\varphi}\right)  \right)  ^{-\frac{1}{2}}\left\vert
T_{1}\varphi_{1}\cdots T_{N}\varphi_{N}\right\rangle
\end{equation}
and
\begin{equation}
\left\vert T_{1}\varphi_{1}\cdots T_{N}\varphi_{N}\right\rangle =\det\left(
\mathbf{\Delta}\left(  \mathbf{\varphi}\right)  \right)  ^{\frac{1}{2}%
}\left\vert \psi_{1}\cdots\psi_{N}\right\rangle
\end{equation}
Hence the action of the $N$-particle diagonal transformation $T_{1}%
\wedge\cdots\wedge T_{N}$ on the IPS\ $\left\vert \varphi_{1}\cdots\varphi
_{N}\right\rangle $ produces the IPS $\left\vert T_{1}\varphi_{1}\cdots
T_{N}\varphi_{N}\right\rangle $ that can also be generated by the action of
the one particle transformation $\mathcal{R}\left(  \mathbf{\varphi}\right)  $
on $\left\vert \varphi_{1}\cdots\varphi_{N}\right\rangle $ where
$\mathcal{R}\left(  \mathbf{\varphi}\right)  =$ $\mathbf{\Delta}\left(
\mathbf{\varphi}\right)  ^{\frac{1}{2}}\mathcal{T}\left(  \mathbf{\varphi
}\right)  $ i.e.
\begin{align}
T_{1}\wedge\cdots\wedge T_{N}\left\vert \varphi_{1}\cdots\varphi
_{N}\right\rangle  &  =\mathcal{R}\left(  \mathbf{\varphi}\right)
\wedge\cdots\wedge\mathcal{R}\left(  \mathbf{\varphi}\right)  \left\vert
\varphi_{1}\cdots\varphi_{N}\right\rangle \nonumber\\
&  =\left\vert \mathcal{R}\left(  \mathbf{\varphi}\right)  \varphi_{1}%
\cdots\mathcal{R}\left(  \mathbf{\varphi}\right)  \varphi_{N}\right\rangle
=\left\vert \mathbf{\Delta}\left(  \mathbf{\varphi}\right)  ^{\frac{1}{2}%
}\mathcal{T}\left(  \mathbf{\varphi}\right)  \varphi_{1}\cdots\mathbf{\Delta
}\left(  \mathbf{\varphi}\right)  ^{\frac{1}{2}}\mathcal{T}\left(
\mathbf{\varphi}\right)  \varphi_{N}\right\rangle \nonumber\\
&  =\det\left(  \mathbf{\Delta}\left(  \mathbf{\varphi}\right)  \right)
^{\frac{1}{2}}\left\vert \mathcal{T}\left(  \mathbf{\varphi}\right)
\varphi_{1}\cdots\mathcal{T}\left(  \mathbf{\varphi}\right)  \varphi
_{N}\right\rangle =\left\vert T_{1}\varphi_{1}\cdots T_{N}\varphi
_{N}\right\rangle
\end{align}
The transformations $\left\{  T_{j}\right\}  $ and $\left\{  T_{1}\wedge
\cdots\wedge T_{N}\right\}  $ form commuting sets, are diagonal in the
space-spin coordinate representation and are semigroups, while the
transformations $\left\{  \mathcal{R}\left(  \mathbf{\varphi}\right)
\right\}  $ and $\left\{  \mathcal{R}\left(  \mathbf{\varphi}\right)
\wedge\cdots\wedge\mathcal{R}\left(  \mathbf{\varphi}\right)  \right\}  $ are
not in general diagonal in this representation nor do they form commuting
sets, but they do form a group.

The correspondence between the $N$-particle transformation $T_{1}\wedge
\cdots\wedge T_{N}$ and the one particle transformation $\mathcal{R}\left(
\mathbf{\varphi}\right)  $ depends on the existence of the $one$ dimensional
representation $\det\left(  \mathbf{\Delta}\right)  $ of $\Delta$ in
$\mathcal{H}^{N}$. This does \emph{not }carry over to the analogous
construction applied to transformations defined in $\mathcal{H}^{2}$ e.g. as
in example 3 of section \ref{Approx}, thus one cannot obtain such a
correspondence between a transformation on $\mathcal{H}^{2}$ and a product of
diagonal transformation acting in $\mathcal{H}^{N}$.

It is interesting to note the form of the first order reduced density
operators, $L_{N}^{1}\left(  D^{N}\right)  $ associated with the normalized
states
\begin{equation}
D^{N}=\frac{\left\vert T_{1}\varphi_{1}\cdots T_{N}\varphi_{N}\right\rangle
\left\langle T_{1}\varphi_{1}\cdots T_{N}\varphi_{N}\right\vert }{\left\langle
T_{1}\varphi_{1}\cdots T_{N}\varphi_{N}|T_{1}\varphi_{1}\cdots T_{N}%
\varphi_{N}\right\rangle }%
\end{equation}
are
\begin{equation}
L_{N}^{1}\left(  D^{N}\right)  =%
%TCIMACRO{\tsum \limits_{1\leq j,k\leq N}}%
%BeginExpansion
{\textstyle\sum\limits_{1\leq j,k\leq N}}
%EndExpansion
\left\vert T_{j}\varphi_{j}\right\rangle \left(  \Delta^{-1}\right)
_{jk}\left\langle T_{k}\varphi_{k}\right\vert =%
%TCIMACRO{\tsum \limits_{1\leq j\leq N}}%
%BeginExpansion
{\textstyle\sum\limits_{1\leq j\leq N}}
%EndExpansion
\left\vert \psi_{j}\right\rangle \left\langle \psi_{j}\right\vert
\label{FORDO}%
\end{equation}
which leads to vastly simplified expressions for expectation values wrt the
state $\left\vert T_{1}\varphi_{1}\cdots T_{N}\varphi_{N}\right\rangle $ than
the general ones developed in section \ref{Example}. Eq. $\left(
\ref{FORDO}\right)  $ covers the case $T_{j}=I_{j}$ $1\leq j\leq N$ leading
to
\begin{equation}
L_{N}^{1}\left(  \frac{\left\vert \varphi_{1}\cdots\varphi_{N}\right\rangle
\left\langle \varphi_{1}\cdots\varphi_{N}\right\vert }{\left\langle
\varphi_{1}\cdots\varphi_{N}|\varphi_{1}\cdots\varphi_{N}\right\rangle
}\right)  =%
%TCIMACRO{\tsum \limits_{1\leq j,k\leq N}}%
%BeginExpansion
{\textstyle\sum\limits_{1\leq j,k\leq N}}
%EndExpansion
\left\vert \varphi_{j}\right\rangle \left(  \Delta^{-1}\right)  _{jk}%
\left\langle \varphi_{k}\right\vert =%
%TCIMACRO{\tsum \limits_{1\leq j\leq N}}%
%BeginExpansion
{\textstyle\sum\limits_{1\leq j\leq N}}
%EndExpansion
\left\vert \psi_{j}\right\rangle \left\langle \psi_{j}\right\vert
\end{equation}
where $\Delta_{jk}=\left\langle \varphi_{j}|\varphi_{k}\right\rangle .$

\section{Examples\label{Example}}

We first consider the approximation indexed by $\left\{  1,\cdots,1\right\}  $
and develop it in manner that can be used for higher order approximations so
that we can see the general structure common to all levels of approximation.
It is clear however that the resulting expressions are more complicated than
the ones that would be obtained if we utilized the correspondence described in
section \ref{NVCT}.

\subsection{Approximations indexed by $\left\{  1,\cdots,1\right\}  $}

The transformations that we consider in this case are
\begin{equation}
\breve{T}_{1}\wedge\cdots\wedge\breve{T}_{N}\left\vert \varphi_{1}%
\cdots\varphi_{N}\right\rangle =\left\vert \breve{T}_{1}\varphi_{1}%
\cdots\breve{T}_{N}\varphi_{N}\right\rangle
\end{equation}
That have integral kernels
\begin{align}
T\left(  \mathsf{z}^{N}\right)   &  =T_{1}\left(  \mathsf{z}_{1}\right)
\vee\cdots\vee T_{N}\left(  \mathsf{z}_{N}\right)  =%
%TCIMACRO{\tsum \limits_{\sigma\in\mathcal{S}_{N}}}%
%BeginExpansion
{\textstyle\sum\limits_{\sigma\in\mathcal{S}_{N}}}
%EndExpansion
T_{1}\left(  \mathsf{z}_{\sigma\left(  1\right)  }\right)  \cdots T_{N}\left(
\mathsf{z}_{\sigma\left(  N\right)  }\right) \nonumber\\
&  =%
%TCIMACRO{\tsum \limits_{\sigma\in\mathcal{S}_{N}}}%
%BeginExpansion
{\textstyle\sum\limits_{\sigma\in\mathcal{S}_{N}}}
%EndExpansion
T_{\sigma\left(  1\right)  }\left(  \mathsf{z}_{1}\right)  \cdots
T_{\sigma\left(  N\right)  }\left(  \mathsf{z}_{N}\right) \nonumber\\
&  =\sum_{\sigma\in\mathcal{S}_{N}}\nu_{\sigma\left(  1\right)  }\left(
\mathsf{z}_{1}\right)  \cdots\nu_{\sigma\left(  N\right)  }\left(
\mathsf{z}_{N}\right)  e^{i\left\{  \omega_{\sigma\left(  1\right)  }\left(
\mathsf{z}_{1}\right)  +\cdots+\omega_{\sigma\left(  N\right)  }\left(
\mathsf{z}_{N}\right)  \right\}  }%
\end{align}
where
\begin{equation}
\omega_{\sigma\left(  j\right)  }\left(  \mathsf{z}_{j}\right)  =\delta\left(
\mathsf{z}_{1}\right)  \cdots\delta\left(  \mathsf{z}_{j-1}\right)
\omega_{\sigma\left(  j\right)  }\left(  \mathsf{z}_{j}\right)  \delta\left(
\mathsf{z}_{j+1}\right)  \cdots\delta\left(  \mathsf{z}_{N}\right)
\end{equation}
and the kernels $\left\{  \nu_{j}\left(  \mathsf{z}_{j}\right)  |1\leq j\leq
N\right\}  $ could be equal to each other or be delta functions as could the
kernels $\left\{  \omega_{j}\left(  \mathsf{z}_{j}\right)  |1\leq j\leq
N\right\}  $. The $N$-particle kernels $\nu,\omega$ are both functions of the
$1$-particle kernels $\left\{  \nu_{1},\ldots,\nu_{N},\omega_{1},\ldots
,\omega_{N}\right\}  $ i.e.%
\begin{align}
\nu &  =\nu\left(  \nu_{1},\ldots,\nu_{N},\omega_{1},\ldots,\omega_{N}\right)
\nonumber\\
\omega &  =\omega\left(  \nu_{1},\ldots,\nu_{N},\omega_{1},\ldots,\omega
_{N}\right)
\end{align}


The first term we need to evaluate is the normalization
\begin{align}
&  Tr\left\{  \mathcal{\breve{D}}^{N}\left(  \nu,\omega\right)  \right\}
=\left\langle \breve{T}\left(  \nu,\omega\right)  \varphi_{1}\cdots\varphi
_{N}\right\vert \left.  \breve{T}\left(  \nu,\omega\right)  \varphi_{1}%
\cdots\varphi_{N}\right\rangle \nonumber\\
&  =\sum_{\sigma,\sigma^{\prime}\in\mathcal{S}_{N}}\left\langle \breve
{T}_{\sigma\left(  1\right)  }\cdots\breve{T}_{\sigma\left(  N\right)
}\varphi_{1}\cdots\varphi_{N}\right.  \left\vert \breve{T}_{\sigma^{\prime
}\left(  1\right)  }\cdots\breve{T}_{\sigma^{\prime}\left(  N\right)  }%
\varphi_{1}\cdots\varphi_{N}\right\rangle \nonumber\\
&  =\sum_{\sigma,\sigma^{\prime}\in\mathcal{S}_{N}}\left\langle \breve
{T}_{\sigma\left(  1\right)  }\varphi_{1}\cdots\breve{T}_{\sigma\left(
N\right)  }\varphi_{N}\right.  \left\vert \breve{T}_{\sigma^{\prime}\left(
1\right)  }\varphi_{1}\cdots\breve{T}_{\sigma^{\prime}\left(  N\right)
}\varphi_{N}\right\rangle \nonumber\\
&  =\left\langle \varphi_{1}\cdots\varphi_{N}\right\vert \left.  \breve
{S}_{eff}\left(  \nu,\omega\right)  \varphi_{1}\cdots\varphi_{N}\right\rangle
\equiv\mathcal{S}\left(  \nu,\omega\right)
\end{align}
where $\left\vert \varphi_{1}\cdots\varphi_{N}\right\rangle $ is a fixed IPS.
The effect of the multiparticle potential $T_{\sigma\left(  1\right)  }\cdots
T_{\sigma\left(  N\right)  }$ on an IPS\ is to produce a new IPS so that we
obtain the inner product between two different IPSs leading to the
normalization
\begin{equation}
\mathcal{S}\left(  \nu,\omega\right)  =%
%TCIMACRO{\tsum \limits_{\sigma,\sigma^{\prime}\in\mathcal{S}_{N}}}%
%BeginExpansion
{\textstyle\sum\limits_{\sigma,\sigma^{\prime}\in\mathcal{S}_{N}}}
%EndExpansion
\det\left\{  \mathbf{S}_{\sigma\sigma^{\prime}}\right\}
\end{equation}
where
\begin{align}
\left(  \mathbf{S}_{\sigma\sigma^{\prime}}\right)  _{jk}  &  =\left\langle
\left.  \nu_{\sigma\left(  j\right)  }e^{i\omega_{\sigma\left(  j\right)  }%
}\varphi_{j}\right\vert \nu_{\sigma^{\prime}\left(  k\right)  }e^{i\omega
_{\sigma^{\prime}\left(  k\right)  }}\varphi_{k}\right\rangle =\nonumber\\
&  \qquad\qquad\qquad\qquad%
%TCIMACRO{\dint }%
%BeginExpansion
{\displaystyle\int}
%EndExpansion
\varphi_{j}\left(  \mathsf{z}\right)  ^{\ast}\varphi_{k}\left(  \mathsf{z}%
\right)  e^{i\left(  \omega_{\sigma^{\prime}\left(  k\right)  }\left(
\mathsf{z}\right)  -\omega_{\sigma\left(  j\right)  }\left(  \mathsf{z}%
\right)  \right)  }\nu_{\sigma\left(  j\right)  }\left(  \mathsf{z}\right)
\nu_{\sigma^{\prime}\left(  k\right)  }\left(  \mathsf{z}\right)  d\mathsf{z}%
\end{align}


The second term produced by the effective Hamiltonian is the expectation value
$\mathcal{K}\left(  \nu,\omega\right)  $ wrt the IPS\ $\left\vert \varphi
_{1}\cdots\varphi_{N}\right\rangle $
\begin{equation}
\mathcal{K}\left(  \nu,\omega\right)  =\mathcal{S}\left(  \nu,\omega\right)
^{-1}\left\langle \varphi_{1}\cdots\varphi_{N}\right\vert \left.  K\left(
\nu,\omega\right)  \varphi_{1}\cdots\varphi_{N}\right\rangle
\end{equation}
of the effective kinetic energy operator $K\left(  \nu,\omega\right)  $ given
by
\begin{align}
\mathcal{K}\left(  \nu,\omega\right)   &  =\mathcal{S}\left(  \nu
,\omega\right)  ^{-1}%
%TCIMACRO{\dsum \limits_{\sigma,\sigma^{\prime}\in\mathcal{S}_{N}}}%
%BeginExpansion
{\displaystyle\sum\limits_{\sigma,\sigma^{\prime}\in\mathcal{S}_{N}}}
%EndExpansion
\left\{  \left\langle \nu_{\sigma\left(  1\right)  }e^{i\omega_{\sigma\left(
1\right)  }}\varphi_{1}\cdots e^{i\omega_{\sigma\left(  N\right)  }}%
\nu_{\sigma\left(  N\right)  }\varphi_{N}\right.  \right. \nonumber\\
&  \qquad\qquad\left.  \left(  -%
%TCIMACRO{\dsum \limits_{1\leq j\leq N}}%
%BeginExpansion
{\displaystyle\sum\limits_{1\leq j\leq N}}
%EndExpansion
\nabla_{j}^{2}\right)  \left.  e^{i\omega_{\sigma^{\prime}\left(  1\right)  }%
}\nu_{\sigma^{\prime}\left(  1\right)  }\varphi_{1}\cdots e^{i\omega
_{\sigma^{\prime}\left(  N\right)  }}\nu_{\sigma\prime\left(  N\right)
}\varphi_{N}\right\rangle \right\}
\end{align}
which can be evaluated by utilizing the expressions described in Appendix
\ref{Transamp} that have been developed by \cite{Verbeek1990} to give
\begin{equation}
\mathcal{K}\left(  \nu,\omega\right)  =%
%TCIMACRO{\tsum \limits_{\sigma,\sigma^{\prime}\in\mathcal{S}_{N}}}%
%BeginExpansion
{\textstyle\sum\limits_{\sigma,\sigma^{\prime}\in\mathcal{S}_{N}}}
%EndExpansion
\operatorname{Tr}\left\{  \mathbf{K}_{\sigma\sigma^{\prime}}\operatorname{adj}%
\left(  \mathbf{S}_{\sigma\sigma^{\prime}}\right)  \right\}
\end{equation}
where
\begin{equation}
\left(  \mathbf{K}_{\sigma\sigma^{\prime}}\right)  _{jk}=-\int\left\{
\nabla\varphi_{j}\left(  \mathsf{z}\right)  ^{\ast}\bullet\nabla\varphi
_{k}\left(  \mathsf{z}\right)  \nu_{\sigma\left(  j\right)  }\left(
\mathsf{z}\right)  \nu_{\sigma^{\prime}\left(  k\right)  }\left(
\mathsf{z}\right)  e^{i\left\{  \omega_{\sigma^{\prime}\left(  k\right)
}\left(  \mathsf{z}\right)  -\,\omega_{\sigma\left(  j\right)  }\left(
\mathsf{z}\right)  \right\}  }\right\}  d\mathsf{z}%
\end{equation}


The expectation value of the one particle potential term
\begin{align}
&  \mathcal{V}_{1}\left(  \nu,\omega\right)  =\mathcal{S}\left(  \nu
,\omega\right)  ^{-1}\nu_{\sigma\left(  1\right)  }\left\langle e^{i\omega
_{\sigma\left(  1\right)  }}\varphi_{1}\cdots e^{i\omega_{\sigma\left(
N\right)  }}\nu_{\sigma\left(  N\right)  }\varphi_{N}\right\vert
\times\nonumber\\
&  \qquad\qquad\qquad\qquad\qquad\qquad\qquad\left.  V_{1}e^{i\omega
_{\sigma^{\prime}\left(  1\right)  }}\nu_{\sigma^{\prime}\left(  1\right)
}\varphi_{1}\cdots e^{i\omega_{\sigma^{\prime}\left(  N\right)  }}\nu
_{\sigma\prime\left(  N\right)  }\varphi_{N}\right\rangle
\end{align}
can be evaluated to give
\begin{equation}
\mathcal{V}_{1}\left(  \nu,\omega\right)  =\mathcal{S}\left(  \nu
,\omega\right)  ^{-1}%
%TCIMACRO{\tsum \limits_{\sigma,\sigma^{\prime}\in\mathcal{S}_{N}}}%
%BeginExpansion
{\textstyle\sum\limits_{\sigma,\sigma^{\prime}\in\mathcal{S}_{N}}}
%EndExpansion
\operatorname{Tr}\left\{  \mathbf{V}_{1\sigma\sigma^{\prime}}%
\operatorname{adj}\left(  \mathbf{S}_{\sigma\sigma^{\prime}}\right)  \right\}
\end{equation}
where $V_{1}\left(  \mathsf{z}^{N}\right)  =%
%TCIMACRO{\dsum \limits_{1\leq j\leq N}}%
%BeginExpansion
{\displaystyle\sum\limits_{1\leq j\leq N}}
%EndExpansion
V_{1}\left(  \mathsf{z}_{j}\right)  $ and
\begin{equation}
\left(  \mathbf{V}_{1\sigma\sigma^{\prime}}\right)  _{jk}=\int\left\{
V_{1}\left(  \mathsf{z}_{1}\right)  \varphi_{j}\left(  \mathsf{z}_{1}\right)
^{\ast}\varphi_{k}\left(  \mathsf{z}_{1}\right)  \nu_{\sigma\left(  j\right)
}\left(  \mathsf{z}_{1}\right)  \nu_{\sigma^{\prime}\left(  k\right)  }\left(
\mathsf{z}_{1}\right)  e^{i\left\{  \omega_{\sigma^{\prime}\left(  k\right)
}\left(  \mathsf{z}\right)  -\,\omega_{\sigma\left(  j\right)  }\left(
\mathsf{z}\right)  \right\}  }\right\}  d\mathsf{z}_{1}%
\end{equation}


The expectation value $\mathcal{V}_{2}\left(  \nu\right)  $ of the coulomb
potential term is given by
\begin{equation}
\mathcal{V}_{2}\left(  \nu,\omega\right)  =\mathcal{S}\left(  \nu
,\omega\right)  ^{-1}\left\langle \nu_{\sigma\left(  1\right)  }%
e^{i\omega_{\sigma\left(  1\right)  }}\varphi_{1}\cdots e^{i\omega
_{\sigma\left(  N\right)  }}\nu_{\sigma\left(  N\right)  }\varphi
_{N}\left\vert V_{2}\right.  e^{i\omega_{\sigma^{\prime}\left(  1\right)  }%
}\nu_{\sigma^{\prime}\left(  1\right)  }\varphi_{1}\cdots e^{i\omega
_{\sigma^{\prime}\left(  N\right)  }}\nu_{\sigma\prime\left(  N\right)
}\varphi_{N}\right\rangle
\end{equation}
which can be evaluated using the expressions in Appendix $\ref{Transamp}$ as
\begin{equation}
\mathcal{S}\left(  \nu,\omega\right)  ^{-1}%
%TCIMACRO{\tsum \limits_{\sigma,\sigma^{\prime}\in\mathcal{S}_{N}}}%
%BeginExpansion
{\textstyle\sum\limits_{\sigma,\sigma^{\prime}\in\mathcal{S}_{N}}}
%EndExpansion
Tr\left\{  \mathbf{G}_{\sigma,\sigma^{\prime}}\operatorname{adj}^{\left(
2\right)  }\left(  \mathbf{S}_{\sigma\sigma^{\prime}}\right)  \right\}
\end{equation}
where
\begin{align}
\left(  \mathbf{G}_{\sigma,\sigma^{\prime}}\right)  _{jklm}  &  =%
%TCIMACRO{\dint }%
%BeginExpansion
{\displaystyle\int}
%EndExpansion
\left\{  \nu_{\sigma\left(  j\right)  }\left(  \mathsf{z}_{1}\right)
\varphi_{j}\left(  \mathsf{z}_{1}\right)  ^{\ast}\nu_{\sigma\left(  k\right)
}\left(  \mathsf{z}_{1}\right)  \varphi_{k}\left(  \mathsf{z}_{1}\right)
^{\ast}\nu_{l}\left(  \mathsf{z}_{2}\right)  \varphi_{\sigma^{\prime}\left(
l\right)  }\left(  \mathsf{z}_{2}\right)  \nu_{m}\left(  \mathsf{z}%
_{2}\right)  \varphi_{\sigma^{\prime}\left(  m\right)  }\left(  \mathsf{z}%
_{2}\right)  \right. \nonumber\\
&  \qquad\qquad\qquad\qquad\qquad\times\frac{e^{i\left\{  \omega
_{\sigma^{\prime}\left(  l\right)  }\left(  \mathsf{z}_{2}\right)
+\omega_{\sigma^{\prime}\left(  m\right)  }\left(  \mathsf{z}_{2}\right)
-\,\omega_{\sigma\left(  j\right)  }\left(  \mathsf{z}_{1}\right)
-\omega_{\sigma\left(  k\right)  }\left(  \mathsf{z}_{1}\right)  \right\}  }%
}{\left\Vert \mathbf{r}_{1}-\mathbf{r}_{2}\right\Vert }\nonumber\\
&  -\nu_{\sigma\left(  j\right)  }\left(  \mathsf{z}_{2}\right)  \varphi
_{j}\left(  \mathsf{z}_{2}\right)  ^{\ast}\nu_{\sigma\left(  k\right)
}\left(  \mathsf{z}_{2}\right)  \varphi_{k}\left(  \mathsf{z}_{2}\right)
^{\ast}\nu_{l}\left(  \mathsf{z}_{1}\right)  \varphi_{\sigma^{\prime}\left(
l\right)  }\left(  \mathsf{z}_{1}\right)  \nu_{m}\left(  \mathsf{z}%
_{1}\right)  \varphi_{\sigma^{\prime}\left(  m\right)  }\left(  \mathsf{z}%
_{1}\right) \nonumber\\
&  \qquad\qquad\qquad\qquad\qquad\times\left.  \frac{e^{i\left\{
\omega_{\sigma^{\prime}\left(  l\right)  }\left(  \mathsf{z}_{1}\right)
+\omega_{\sigma^{\prime}\left(  m\right)  }\left(  \mathsf{z}_{1}\right)
-\,\omega_{\sigma\left(  j\right)  }\left(  \mathsf{z}_{2}\right)
-\omega_{\sigma\left(  k\right)  }\left(  \mathsf{z}_{2}\right)  \right\}  }%
}{\left\Vert \mathbf{r}_{1}-\mathbf{r}_{2}\right\Vert }\right\}
d\mathsf{z}_{1}d\mathsf{z}_{2}%
\end{align}
The value of the energy functional $E\left(  \nu,\omega\right)  $ is then the
sum
\begin{equation}
E\left(  \nu,\omega\right)  =H_{0}+\mathcal{K}\left(  \nu,\omega\right)
+\mathcal{V}_{1}\left(  \nu,\omega\right)  +\mathcal{V}_{2}\left(  \nu
,\omega\right)
\end{equation}
where the kernels $\left\{  \nu,\omega\right\}  $ are expressed in terms of
$\left\{  \nu_{1},\ldots,\nu_{N},\omega_{1},\ldots,\omega_{N}\right\}  $ which
in turn can be expanded in terms of a fixed subset of $s$ functions $\left\{
f_{1},\ldots,f_{s}\right\}  $ and the coefficients of these expansions then
become the variational parameters of the problem.

\subsection{Approximations indexed by $\left\{  2,\cdots,2\right\}  $}

The transformations that we consider in this case are
\begin{equation}
\breve{T}_{1}\wedge\cdots\wedge\breve{T}_{\frac{N}{2}}\left\vert \varphi
_{1}\cdots\varphi_{N}\right\rangle =%
%TCIMACRO{\dsum \limits_{\sigma\in\mathcal{S}_{N}}}%
%BeginExpansion
{\displaystyle\sum\limits_{\sigma\in\mathcal{S}_{N}}}
%EndExpansion
\left\vert \breve{T}_{1}\left(  \varphi_{\sigma\left(  1\right)  }%
\varphi_{\sigma\left(  2\right)  }\right)  \cdots\breve{T}_{\frac{N}{2}%
}\left(  \varphi_{\sigma\left(  N-1\right)  }\varphi_{\sigma\left(  N\right)
}\right)  \right\rangle
\end{equation}
where the amplitudes of $\left\vert T_{j}\left(  \varphi_{\sigma\left(
2j-1\right)  }\varphi_{\sigma\left(  2j\right)  }\right)  \right\rangle $ are
given by
\begin{equation}
\left\langle \mathsf{z}_{2j-1},\mathsf{z}_{2j}\right.  \left\vert T_{j}\left(
\varphi_{\sigma\left(  2j-1\right)  }\varphi_{\sigma\left(  2j\right)
}\right)  \right\rangle =\int T_{j}\left(  \mathsf{z}_{2j-1},\mathsf{z}%
_{2j}\right)  \varphi_{\sigma\left(  2j-1\right)  }\left(  \mathsf{z}%
_{2j-1}\right)  \varphi_{\sigma\left(  2j\right)  }\left(  \mathsf{z}%
_{2j}\right)  d\mathsf{z}_{2j-1}d\mathsf{z}_{2j}%
\end{equation}
and the integral kernels of $\breve{T}_{1}\wedge\cdots\wedge\breve{T}%
_{\frac{N}{2}}$ are
\begin{align}
T\left(  \mathsf{z}^{N}\right)   &  =T_{1}\left(  \mathsf{z}_{1}%
,\mathsf{z}_{2}\right)  \vee\cdots\vee T_{\frac{N}{2}}\left(  \mathsf{z}%
_{N-1,}\mathsf{z}_{N}\right)  =%
%TCIMACRO{\tsum \limits_{\sigma\in\mathcal{S}_{N}}}%
%BeginExpansion
{\textstyle\sum\limits_{\sigma\in\mathcal{S}_{N}}}
%EndExpansion
T_{1}\left(  \mathsf{z}_{\sigma\left(  1\right)  },\mathsf{z}_{\sigma\left(
2\right)  }\right)  \cdots T_{\frac{N}{2}}\left(  \mathsf{z}_{\sigma\left(
N-1\right)  ,}\mathsf{z}_{\sigma\left(  N\right)  }\right) \nonumber\\
&  =\sum_{\sigma\in\mathcal{S}_{N}}\nu_{1}\left(  \mathsf{z}_{\sigma\left(
1\right)  },\mathsf{z}_{\sigma\left(  2\right)  }\right)  \cdots\nu_{\frac
{N}{2}}\left(  \mathsf{z}_{\sigma\left(  N-1\right)  },\mathsf{z}%
_{\sigma\left(  2N\right)  }\right)  e^{i\left\{  \omega_{1}\left(
\mathsf{z}_{\sigma\left(  1\right)  },\mathsf{z}_{\sigma\left(  2\right)
}\right)  +\cdots+\omega_{\frac{N}{2}}\left(  \mathsf{z}_{\sigma\left(
N-1\right)  ,}\mathsf{z}_{\sigma\left(  N\right)  }\right)  \right\}  }%
\end{align}
where
\begin{equation}
\omega_{j}\left(  \mathsf{z}_{\sigma\left(  2j-1\right)  },\mathsf{z}%
_{\sigma\left(  2j\right)  }\right)  =\delta\left(  \mathsf{z}_{1}\right)
\cdots\omega_{j}\left(  \mathsf{z}_{\sigma\left(  2j-1\right)  }%
,\mathsf{z}_{\sigma\left(  2j\right)  }\right)  \cdots\delta\left(
\mathsf{z}_{N}\right)
\end{equation}


In this case we must treat the effective overlap operator $S\left(  \nu
,\omega\right)  $ as a full $N$-particle operator
\begin{align}
S\left(  \nu,\omega\right)   &  =\int\left\{
%TCIMACRO{\tsum \limits_{\sigma,\sigma^{\prime}\in\mathcal{S}_{N}}}%
%BeginExpansion
{\textstyle\sum\limits_{\sigma,\sigma^{\prime}\in\mathcal{S}_{N}}}
%EndExpansion
T_{1}\left(  \mathsf{z}_{\sigma\left(  1\right)  },\mathsf{z}_{\sigma\left(
2\right)  }\right)  ^{\ast}\cdots T_{\frac{N}{2}}\left(  \mathsf{z}%
_{\sigma\left(  N-1\right)  ,}\mathsf{z}_{\sigma\left(  N\right)  }\right)
^{\ast}\right. \nonumber\\
&  \left.  T_{1}\left(  \mathsf{z}_{\sigma^{\prime}\left(  1\right)
},\mathsf{z}_{\sigma^{\prime}\left(  2\right)  }\right)  \cdots T_{\frac{N}%
{2}}\left(  \mathsf{z}_{\sigma^{\prime}\left(  N-1\right)  ,}\mathsf{z}%
_{\sigma^{\prime}\left(  N\right)  }\right)  \left\vert \mathsf{z}%
^{N}\right\rangle \left\langle \mathsf{z}^{N}\right\vert \right\}
d\mathsf{z}^{N}%
\end{align}
and take the expectation value of this operator wrt the IPS $\left\vert
\varphi_{1}\cdots\varphi_{N}\right\rangle $ producing
\begin{align}
&  \mathcal{S}\left(  \nu,\omega\right)  =\left\langle \varphi_{1}%
\cdots\varphi_{N}|S\left(  \nu,\omega\right)  \varphi_{1}\cdots\varphi
_{N}\right\rangle =\nonumber\\
&
%TCIMACRO{\tsum \limits_{\sigma,\sigma^{\prime},\sigma^{\prime\prime}%
%\in\mathcal{S}_{N}}}%
%BeginExpansion
{\textstyle\sum\limits_{\sigma,\sigma^{\prime},\sigma^{\prime\prime}%
\in\mathcal{S}_{N}}}
%EndExpansion%
%TCIMACRO{\dint }%
%BeginExpansion
{\displaystyle\int}
%EndExpansion
\left\{  \varphi_{_{1}}\left(  \mathsf{z}_{1}\right)  ^{\ast}\cdots\varphi
_{N}\left(  \mathsf{z}_{N}\right)  ^{\ast}\left(  -1\right)  ^{\pi\left(
\sigma^{\prime\prime}\right)  }\varphi_{1}\left(  \mathsf{z}_{\sigma
^{\prime\prime}\left(  1\right)  }\right)  \cdots\varphi_{N}\left(
\mathsf{z}_{\sigma^{\prime\prime}\left(  N\right)  }\right)  \right.
\nonumber\\
&  \qquad\qquad\qquad\qquad\qquad\qquad\times%
%TCIMACRO{\dprod \limits_{1\leq k\leq\frac{N}{2}}}%
%BeginExpansion
{\displaystyle\prod\limits_{1\leq k\leq\frac{N}{2}}}
%EndExpansion
\nu_{k}\left(  \mathsf{z}_{\sigma\left(  2k-1\right)  },\mathsf{z}%
_{\sigma\left(  2k\right)  }\right)  \nu_{k}\left(  \mathsf{z}_{\sigma
^{\prime}\left(  2k-1\right)  },\mathsf{z}_{\sigma^{\prime}\left(  2k\right)
}\right) \nonumber\\
&  \qquad\times\left.  e^{i\left\{  \omega_{1}\left(  \mathsf{z}%
_{\sigma^{\prime}\left(  1\right)  },\mathsf{z}_{\sigma^{\prime}\left(
2\right)  }\right)  -\omega_{1}\left(  \mathsf{z}_{\sigma\left(  1\right)
},\mathsf{z}_{\sigma\left(  2\right)  }\right)  +\cdots+\omega_{\frac{N}{2}%
}\left(  \mathsf{z}_{\sigma^{\prime}\left(  N-1\right)  ,}\mathsf{z}%
_{\sigma^{\prime}\left(  N\right)  }\right)  -\omega_{\frac{N}{2}}\left(
\mathsf{z}_{\sigma\left(  N-1\right)  ,}\mathsf{z}_{\sigma\left(  N\right)
}\right)  \right\}  }\right\}  d\mathsf{z}^{N}%
\end{align}
The effective kinetic energy is the expectation value $\mathcal{K}\left(
\nu,\omega\right)  $
\begin{equation}
\mathcal{K}\left(  \nu,\omega\right)  =\mathcal{S}\left(  \nu,\omega\right)
^{-1}\left\langle \varphi_{1}\cdots\varphi_{N}\right\vert \left.  K\left(
\nu,\omega\right)  \varphi_{1}\cdots\varphi_{N}\right\rangle
\end{equation}
of the effective kinetic energy operator $K\left(  \nu,\omega\right)  $ given
by%
\begin{align}
K\left(  \nu,\omega\right)   &  =\,\underset{I}{\underbrace{-\nabla
\bullet\left(  \nu^{2}\right)  \mathbb{I}\bullet\nabla}}-\underset
{II}{\underbrace{\nabla\bullet\left(  \nu^{2}\nabla\omega-i\nu\nabla
\nu\right)  }}\nonumber\\
&  -\underset{III}{\underbrace{\left(  \nu^{2}\nabla\omega-i\nu\nabla
\nu\right)  \bullet\nabla}}+\underset{IV}{\underbrace{\left(  \nu\nabla
\omega-i\nabla\nu\right)  \bullet\mathbb{I}\bullet\left(  \nu\nabla
\omega-i\nabla\nu\right)  }} \label{KE2}%
\end{align}
where the $N-$particle integral kernels $\nu,\omega$ are functions of the two
particle integral kernels $\left\{  \nu_{1},\ldots,\nu_{\frac{N}{2}}%
,\omega_{1},\ldots,\omega_{\frac{N}{2}}\right\}  $ i.e.%
\begin{align}
\nu &  =\nu\left(  \nu_{1},\ldots,\nu_{\frac{N}{2}},\omega_{1},\ldots
,\omega_{\frac{N}{2}}\right) \nonumber\\
\omega &  =\omega\left(  \nu_{1},\ldots,\nu_{\frac{N}{2}},\omega_{1}%
,\ldots,\omega_{\frac{N}{2}}\right)
\end{align}
The expansions of the various parts of the effective kinetic energy operator
eq. $\left(  \ref{KE2}\right)  $ are shown in Appendix \ref{KinEn}. The
expectation value of the one particle potential
\begin{equation}
\mathcal{V}_{1}\left(  \nu,\omega\right)  =\mathcal{S}\left(  \nu
,\omega\right)  ^{-1}\left\langle \varphi_{1}\cdots\varphi_{N}\left\vert
V_{1}\right.  \varphi_{1}\cdots\varphi_{N}\right\rangle
\end{equation}
is given in terms of the effective operators%
\begin{align}
&  V_{1}=%
%TCIMACRO{\dsum \limits_{\sigma,\sigma^{\prime}\in\mathcal{S}_{N}}}%
%BeginExpansion
{\displaystyle\sum\limits_{\sigma,\sigma^{\prime}\in\mathcal{S}_{N}}}
%EndExpansion%
%TCIMACRO{\dsum \limits_{1\leq j\leq N}}%
%BeginExpansion
{\displaystyle\sum\limits_{1\leq j\leq N}}
%EndExpansion
\int\left\{  V_{1}\left(  \mathsf{z}_{j}\right)  T_{1}\left(  \mathsf{z}%
_{\sigma\left(  1\right)  },\mathsf{z}_{\sigma\left(  2\right)  }\right)
^{\ast}\cdots T_{\frac{N}{2}}\left(  \mathsf{z}_{\sigma\left(  N-1\right)
,}\mathsf{z}_{\sigma\left(  N\right)  }\right)  ^{\ast}\right. \nonumber\\
&  \qquad\qquad\qquad\qquad\left.  T_{1}\left(  \mathsf{z}_{\sigma^{\prime
}\left(  1\right)  },\mathsf{z}_{\sigma^{\prime}\left(  2\right)  }\right)
\cdots T_{\frac{N}{2}}\left(  \mathsf{z}_{\sigma^{\prime}\left(  N-1\right)
,}\mathsf{z}_{\sigma^{\prime}\left(  N\right)  }\right)  \left\vert
\mathsf{z}^{N}\right\rangle \left\langle \mathsf{z}^{N}\right\vert \right\}
d\mathsf{z}^{N}%
\end{align}
and can be evaluated to give
\begin{align}
\mathcal{V}_{1}\left(  \nu,\omega\right)   &  =\mathcal{S}\left(  \nu
,\omega\right)  ^{-1}\times\nonumber\\
&
%TCIMACRO{\dsum \limits_{\sigma,\sigma^{\prime},\sigma^{\prime\prime}%
%\in\mathcal{S}_{N}}}%
%BeginExpansion
{\displaystyle\sum\limits_{\sigma,\sigma^{\prime},\sigma^{\prime\prime}%
\in\mathcal{S}_{N}}}
%EndExpansion%
%TCIMACRO{\dsum \limits_{1\leq j\leq N}}%
%BeginExpansion
{\displaystyle\sum\limits_{1\leq j\leq N}}
%EndExpansion
\int\left\{  \varphi_{_{1}}\left(  \mathsf{z}_{1}\right)  ^{\ast}\cdots
\varphi_{N}\left(  \mathsf{z}_{N}\right)  ^{\ast}\left(  -1\right)
^{\pi\left(  \sigma^{\prime\prime}\right)  }\varphi_{1}\left(  \mathsf{z}%
_{\sigma^{\prime\prime}\left(  1\right)  }\right)  \cdots\varphi_{N}\left(
\mathsf{z}_{\sigma^{\prime\prime}\left(  N\right)  }\right)  \right.
\nonumber\\
&  \qquad\times V_{1}\left(  \mathsf{z}_{j}\right)
%TCIMACRO{\dprod \limits_{1\leq k\leq\frac{N}{2}}}%
%BeginExpansion
{\displaystyle\prod\limits_{1\leq k\leq\frac{N}{2}}}
%EndExpansion
\nu_{k}\left(  \mathsf{z}_{\sigma\left(  2k-1\right)  },\mathsf{z}%
_{\sigma\left(  2k\right)  }\right)  \nu_{k}\left(  \mathsf{z}_{\sigma
^{\prime}\left(  2k-1\right)  },\mathsf{z}_{\sigma^{\prime}\left(  2k\right)
}\right) \nonumber\\
&  \qquad\left.  e^{i\left\{  \omega_{1}\left(  \mathsf{z}_{\sigma^{\prime
}\left(  1\right)  },\mathsf{z}_{\sigma^{\prime}\left(  2\right)  }\right)
-\omega_{1}\left(  \mathsf{z}_{\sigma\left(  1\right)  },\mathsf{z}%
_{\sigma\left(  2\right)  }\right)  +\cdots+\omega_{\frac{N}{2}}\left(
\mathsf{z}_{\sigma^{\prime}\left(  N-1\right)  ,}\mathsf{z}_{\sigma^{\prime
}\left(  N\right)  }\right)  -\omega_{\frac{N}{2}}\left(  \mathsf{z}%
_{\sigma\left(  N-1\right)  ,}\mathsf{z}_{\sigma\left(  N\right)  }\right)
\right\}  }\right\}  d\mathsf{z}^{N}%
\end{align}


The expectation value $\mathcal{V}_{2}\left(  \nu\right)  $ of the coulomb
potential term is given by
\begin{equation}
\mathcal{V}_{2}\left(  \nu,\omega\right)  =\mathcal{S}\left(  \nu
,\omega\right)  ^{-1}\left\langle \varphi_{1}\cdots\varphi_{N}\left\vert
V_{2}\right.  \varphi_{1}\cdots\varphi_{N}\right\rangle
\end{equation}
which can be evaluated in terms of the effective coulomb operator
\begin{align}
&
%TCIMACRO{\dsum \limits_{\sigma,\sigma^{\prime}\in\mathcal{S}_{N}}}%
%BeginExpansion
{\displaystyle\sum\limits_{\sigma,\sigma^{\prime}\in\mathcal{S}_{N}}}
%EndExpansion%
%TCIMACRO{\dsum \limits_{1\leq j<k\leq N}}%
%BeginExpansion
{\displaystyle\sum\limits_{1\leq j<k\leq N}}
%EndExpansion
\int\left\{  V_{2}\left(  \mathsf{z}_{j,}\mathsf{z}_{k}\right)  T_{1}\left(
\mathsf{z}_{\sigma\left(  1\right)  },\mathsf{z}_{\sigma\left(  2\right)
}\right)  ^{\ast}\cdots T_{\frac{N}{2}}\left(  \mathsf{z}_{\sigma\left(
N-1\right)  ,}\mathsf{z}_{\sigma\left(  N\right)  }\right)  ^{\ast}\right.
\nonumber\\
&  \qquad\qquad\qquad\qquad\left.  T_{1}\left(  \mathsf{z}_{\sigma^{\prime
}\left(  1\right)  },\mathsf{z}_{\sigma^{\prime}\left(  2\right)  }\right)
\cdots T_{\frac{N}{2}}\left(  \mathsf{z}_{\sigma^{\prime}\left(  N-1\right)
,}\mathsf{z}_{\sigma^{\prime}\left(  N\right)  }\right)  \left\vert
\mathsf{z}^{N}\right\rangle \left\langle \mathsf{z}^{N}\right\vert \right\}
d\mathsf{z}^{N}%
\end{align}
to give%
\begin{align}
&  \mathcal{V}_{2}\left(  \nu,\omega\right)  =\mathcal{S}\left(  \nu
,\omega\right)  ^{-1}\times\nonumber\\
&
%TCIMACRO{\dsum \limits_{\sigma,\sigma^{\prime},\sigma^{\prime\prime}%
%\in\mathcal{S}_{N}}}%
%BeginExpansion
{\displaystyle\sum\limits_{\sigma,\sigma^{\prime},\sigma^{\prime\prime}%
\in\mathcal{S}_{N}}}
%EndExpansion%
%TCIMACRO{\dsum \limits_{1\leq j<k\leq N}}%
%BeginExpansion
{\displaystyle\sum\limits_{1\leq j<k\leq N}}
%EndExpansion
\int\left\{  \varphi_{_{1}}\left(  \mathsf{z}_{1}\right)  ^{\ast}\cdots
\varphi_{N}\left(  \mathsf{z}_{N}\right)  ^{\ast}\left(  -1\right)
^{\pi\left(  \sigma^{\prime\prime}\right)  }\varphi_{1}\left(  \mathsf{z}%
_{\sigma^{\prime\prime}\left(  1\right)  }\right)  \cdots\varphi_{N}\left(
\mathsf{z}_{\sigma^{\prime\prime}\left(  N\right)  }\right)  \right.
\nonumber\\
&  \qquad\qquad\qquad\qquad\qquad\times V_{2}\left(  \mathsf{z}_{j}%
,\mathsf{z}_{k}\right)
%TCIMACRO{\dprod \limits_{1\leq k\leq\frac{N}{2}}}%
%BeginExpansion
{\displaystyle\prod\limits_{1\leq k\leq\frac{N}{2}}}
%EndExpansion
\nu_{k}\left(  \mathsf{z}_{\sigma\left(  2k-1\right)  },\mathsf{z}%
_{\sigma\left(  2k\right)  }\right)  \nu_{k}\left(  \mathsf{z}_{\sigma
^{\prime}\left(  2k-1\right)  },\mathsf{z}_{\sigma^{\prime}\left(  2k\right)
}\right) \nonumber\\
&  \qquad\qquad\qquad\left.  e^{i\left\{  \omega_{1}\left(  \mathsf{z}%
_{\sigma^{\prime}\left(  1\right)  },\mathsf{z}_{\sigma^{\prime}\left(
2\right)  }\right)  -\omega_{1}\left(  \mathsf{z}_{\sigma\left(  1\right)
},\mathsf{z}_{\sigma\left(  2\right)  }\right)  +\cdots+\omega_{\frac{N}{2}%
}\left(  \mathsf{z}_{\sigma^{\prime}\left(  N-1\right)  ,}\mathsf{z}%
_{\sigma^{\prime}\left(  N\right)  }\right)  -\omega_{\frac{N}{2}}\left(
\mathsf{z}_{\sigma\left(  N-1\right)  ,}\mathsf{z}_{\sigma\left(  N\right)
}\right)  \right\}  }\right\}  d\mathsf{z}^{N}%
\end{align}
where%
\begin{equation}
V_{2}\left(  \mathsf{z}_{j},\mathsf{z}_{k}\right)  =\frac{1}{\left\Vert
\mathbf{r}_{j}-\mathbf{r}_{k}\right\Vert }\delta\left(  \xi_{j}-\xi
_{k}\right)
\end{equation}
for $\mathsf{z}_{j}\equiv\left(  \mathbf{r}_{j},\xi_{j}\right)  $ and
$\mathsf{z}_{k}\equiv\left(  \mathbf{r}_{k},\xi_{k}\right)  .$

The value of the energy functional $E\left(  \nu,\omega\right)  $ is then the
sum
\begin{equation}
E\left(  \mathcal{\Psi}\left(  \nu_{1},\ldots,\nu_{\frac{N}{2}},\omega
_{1},\ldots,\omega_{\frac{N}{2}}\right)  \right)  =H_{0}+\mathcal{K}\left(
\nu,\omega\right)  +\mathcal{V}_{1}\left(  \nu,\omega\right)  +\mathcal{V}%
_{2}\left(  \nu,\omega\right)
\end{equation}
and the optimal IPS\ is given by minimizing $E\left(  \mathcal{\Psi}\left(
\nu,\omega\right)  \right)  $ over $\nu=\left(  \nu_{1},\cdots,\nu\frac{_{N}%
}{2}\right)  $ and $\omega=\left(  \omega_{1},\cdots,\omega\frac{_{N}}%
{2}\right)  \mathsf{.}$ In an analogous fashion to the $\left\{
1,\ldots,1\right\}  $ case the functions $\left\{  \nu_{1},\cdots,\nu
\frac{_{N}}{2},\omega_{1},\cdots,\omega\frac{_{N}}{2}\right\}  $ can be
expanded in terms of a fixed set of $s$ functions $\left\{  f_{1},\ldots
,f_{s}\right\}  ,$ where now these functions are two particle functions.

\section{Summary\label{Sum}}

In this article we have shown that

\begin{description}
\item[$\bullet$] one can generate all pure states from a cyclic reference
state (a pure state that has a wavefunction that is non zero almost
everywhere) by transformations that are diagonal in the space-spin coordinate representation,

\item[$\bullet$] these transformations, which form an abelian semigroup
$C_{c}$, also maintain the local form of the defining space-spin coordinate
representation of the Hamiltonian, leading to effective Hamiltonians that have
direct physical interpretability in terms of momentum and potential operators
and thus have a certain classical flavor,

\item[$\bullet$] one can form an approximation theory by restricting attention
to sub semigroups of $C_{c}$, that lead to approximations which have a
\emph{direct} physical interpretation.
\end{description}

These properties should be contrasted to the effect of the transformation of
state in truncated configurational interaction, Multi Configurational Self
Consistent Field and couple cluster methods that do not lead to or matrix
elements that could identified with such effective Hamiltonians.

To illuminate these points we presented some examples that showed the form of
terms that need to be calculated in order to produce an energy functional that
can be optimized. Many strategies are obviously possible for the
simplification of the expressions we obtained which were, by intention, in a
very raw form that need to be simplified in order to produce practical
computational formulae.

The approximations that are suggested in this article present a new approach
to obtaining approximate solutions to the Schr\"{o}dinger equation, and offer
possibilities for development in many directions.

\section{Acknowledgements}

The author would like to thank Dr Sam Trickey for useful discussions
pertaining to the ideas presented in this article.

\appendix

\section{Subspaces of Operators\label{spaces}}

The set of trace class operators is defined as
\begin{equation}
\mathcal{B}_{1}\left(  \mathcal{H}\right)  =\left\{  X|X\in\mathcal{B}\left(
\mathcal{H}\right)  ;Tr\left\{  \left(  X^{\dagger}X\right)  ^{\frac{1}{2}%
}\right\}  <\infty\right\}
\end{equation}
where $\mathcal{B}\left(  \mathcal{H}\right)  $ is the space of bounded
operators acting in the Hilbert space $\mathcal{H}$. $\mathcal{B}_{1}\left(
\mathcal{H}\right)  $ is a Banach space equipped with the norm
\begin{equation}
\left\Vert X\right\Vert _{1}=Tr\left\{  \left(  X^{\dagger}X\right)
^{\frac{1}{2}}\right\}
\end{equation}
The space $\mathcal{B}\left(  \mathcal{H}\right)  $ is also a Banach space
with a norm defined as
\begin{equation}
\left\Vert X\right\Vert =\sup_{\left\vert \Psi\right\rangle \in\mathcal{H}%
}\left\{  \frac{\left\langle \Psi|X\Psi\right\rangle }{\left\langle \Psi
|\Psi\right\rangle }\right\}
\end{equation}
The closure of $\mathcal{B}_{1}\left(  \mathcal{H}\right)  $ wrt $\left\Vert
{}\right\Vert $ produces $\mathcal{B}\left(  \mathcal{H}\right)  $. The
relationship between these spaces is
\begin{equation}
\mathcal{B}\left(  \mathcal{H}\right)  \supseteq\mathcal{B}_{1}\left(
\mathcal{H}\right)
\end{equation}
and their norms is
\begin{equation}
\left\Vert X\right\Vert \leq\left\Vert X\right\Vert _{1}%
\end{equation}


\section{States\label{states}}

The states of an $N$-particle quantum mechanical system correspond to
positive, normalized trace class operators and form a convex subset,
$\mathcal{S}_{N}$, of the Banach space, $\mathcal{B}_{1}\left(  \mathcal{H}%
^{N}\right)  $, of $N$-particle trace class operators, defined as%
\begin{equation}
\mathcal{S}_{N}=\left\{  D^{N}|D^{N}\in\mathcal{B}_{1}\left(  \mathcal{H}%
^{N}\right)  ,D^{N}\geq0,TrD^{N}=1\right\}
\end{equation}
Pure states are idempotent operators i.e.
\begin{equation}
\left(  D^{N}\right)  ^{2}=D^{N}%
\end{equation}
and \emph{normalized }pure states\emph{ }can be associated with, in a 1-1
fashion, with a ray $\left\{  \alpha\left\vert \Psi\right\rangle |\alpha
\in\mathbb{C}\right\}  $ generated by a vector $\left\vert \Psi\right\rangle $
in the N-particle Hilbert space $\mathcal{H}^{N}$ by the relationship
\begin{equation}
D^{N}=\frac{\left\vert \Psi\right\rangle \left\langle \Psi\right\vert
}{\left\langle \Psi|\Psi\right\rangle }%
\end{equation}
The set of \textit{unnormalized }states form the set of positive trace class
operators, $\mathcal{B}_{1+}\left(  \mathcal{H}^{1}\right)  $. States that are
not idempotent are called mixed or ensemble states and correspond to
statistical mixtures of pure states.

\section{Discrete versus Continuous \label{Difference}}

The discrete semigroup $C_{d}$ is given by operators of the form
\begin{equation}
Q=%
%TCIMACRO{\dsum \limits_{1\leq j\leq\infty}}%
%BeginExpansion
{\displaystyle\sum\limits_{1\leq j\leq\infty}}
%EndExpansion
q_{j}\left\vert e_{j}\right\rangle \left\langle e_{j}\right\vert
\end{equation}
and the continuous $C_{c}$ by
\begin{equation}
T=%
%TCIMACRO{\dint }%
%BeginExpansion
{\displaystyle\int}
%EndExpansion
T\left(  x\right)  \left\vert x\right\rangle \left\langle x\right\vert dx
\end{equation}
In the discrete case an invariant one dimensional subspace exists if there is
a vector $\left\vert v\right\rangle $ such that%
\begin{equation}
Q\left\vert v\right\rangle =\alpha_{Q}\left\vert v\right\rangle \qquad\forall
Q\in C_{d}%
\end{equation}
This condition can be expressed as
\begin{equation}%
%TCIMACRO{\dsum \limits_{1\leq j\leq\infty}}%
%BeginExpansion
{\displaystyle\sum\limits_{1\leq j\leq\infty}}
%EndExpansion
Q_{j}\left\vert e_{j}\right\rangle \left\langle e_{j}\right\vert \left.
v\right\rangle =\alpha_{Q}%
%TCIMACRO{\dsum \limits_{1\leq j\leq\infty}}%
%BeginExpansion
{\displaystyle\sum\limits_{1\leq j\leq\infty}}
%EndExpansion
\left\vert e_{j}\right\rangle \left\langle e_{j}\right\vert \left.
v\right\rangle \qquad\forall Q\in C_{d}%
\end{equation}
as the basis $\left\{  e_{j}\right\}  $ is linearly independent we can treat
each component separately to obtain%
\begin{equation}
Q_{j}\left\vert e_{j}\right\rangle \left\langle e_{j}\right\vert \left.
v\right\rangle =\alpha_{Q}\left\vert e_{j}\right\rangle \left\langle
e_{j}\right\vert \left.  v\right\rangle \qquad\forall Q_{j}\quad\forall j
\end{equation}
this implies that
\begin{equation}
Q_{j}=\alpha_{Q}\qquad\forall v_{j}\neq0\quad\forall Q_{j} \label{q}%
\end{equation}
let
\begin{equation}
v_{1}\neq0\quad\&\quad v_{2}\neq0
\end{equation}
and all other expansion coefficients of $\left\vert v\right\rangle $ equal to
zero then eq. $\left(  \ref{q}\right)  $
\begin{equation}
Q_{1}=Q_{2}=\alpha_{Q}\quad\forall Q_{1},Q_{2} \label{multi}%
\end{equation}
which clearly is not possible. Hence only one coefficient can be nonzero and
we obtain%
\begin{equation}
\left\vert v\right\rangle =\beta\left\vert e_{k}\right\rangle
\end{equation}
giving
\begin{equation}%
%TCIMACRO{\dsum \limits_{1\leq j\leq\infty}}%
%BeginExpansion
{\displaystyle\sum\limits_{1\leq j\leq\infty}}
%EndExpansion
Q_{j}\left\vert e_{j}\right\rangle \left\langle e_{j}\right\vert \beta\left.
e_{k}\right\rangle =Q_{k}\beta\left\vert e_{k}\right\rangle \quad\forall Q\in
C_{d}%
\end{equation}
which shows that $\alpha_{Q}=Q_{k}$. At the opposite extreme if all the
expansion coefficients $\left\{  v_{j}\right\}  $ are non zero then
$\left\vert v\right\rangle $ would be a cyclic vector for $C_{d}$ i.e. the
Hilbert space is generated by the action of $C_{d}$ on $\left\vert
v\right\rangle $.

In the continuous case an invariant one dimensional subspace exists if
\begin{equation}
T\left\vert v\right\rangle =\alpha_{T}\left\vert v\right\rangle \qquad\forall
T\in C_{c}%
\end{equation}
which can be expressed as
\begin{align}%
%TCIMACRO{\dint }%
%BeginExpansion
{\displaystyle\int}
%EndExpansion
T\left(  x\right)  \left\vert x\right\rangle \left\langle x\right\vert \left.
v\right\rangle dx  &  =\alpha_{T}%
%TCIMACRO{\dint }%
%BeginExpansion
{\displaystyle\int}
%EndExpansion
\left\vert x\right\rangle \left\langle x\right\vert \left.  v\right\rangle
dx\qquad\forall T\in C_{c}\\
&  \Rightarrow t\left(  x\right)  =\alpha_{T}\quad v\left(  x\right)
\neq0\quad\forall t,\quad x\in\Delta
\end{align}
so that for sets of non zero measure $\Delta$
\begin{equation}
T\left(  x\right)  \left\langle x\right\vert \left.  v\right\rangle
=\alpha_{T}\left\langle x\right\vert \left.  v\right\rangle \quad\forall
T,\quad x\in\Delta
\end{equation}
and
\begin{equation}
T\left(  x\right)  =\alpha_{T}\quad\text{when }v\left(  x\right)  \neq
0\quad,x\in\Delta\quad\forall T
\end{equation}
However now there no single components $\left\{  \left\vert x\right\rangle
\right\}  $ as these vectors do not belong to the Hilbert space and there is
no such square integrable function $v\left(  x\right)  $ that makes this true
for all $T$ (the analogous condition to eq. $\left(  \ref{multi}\right)  $ in
the continuous case cannot be satisfied), so no invariant one dimensional
Hilbert subspaces for the semigroup $\left\{  T\right\}  $ exist. If the
square integrable function $v\left(  x\right)  $ is non zero almost everywhere
then it is a cyclic vector for the semigroup $C_{c}.$

\section{Transformation of the Kinetic Energy Operator\label{KinEnTrans}}

\subsection{Transition Amplitudes}

We consider the commutation relationships between the operators $\mathbf{P,}%
e^{i\breve{\omega}_{N}}$ and $\breve{\nu}_{N}$ and observe that

\begin{enumerate}
\item
\begin{align}
&  \mathbf{P}e^{i\breve{\omega}_{N}}\left\vert \Psi\right\rangle =\left(
-i\nabla\right)  e^{i\breve{\omega}_{N}}\left\vert \Psi\right\rangle
=e^{i\breve{\omega}_{N}}\left(  -i\nabla\right)  \left\vert \Psi\right\rangle
+e^{i\breve{\omega}_{N}}\left(  \nabla\breve{\omega}_{N}\right)  \left\vert
\Psi\right\rangle \nonumber\\
&  \qquad\qquad\qquad\left.  =\right.  e^{i\breve{\omega}_{N}}\left(
-i\nabla+\nabla\breve{\omega}_{N}\right)  \left\vert \Psi\right\rangle
=e^{i\breve{\omega}_{N}}\left\{  \mathbf{P}+\nabla\breve{\omega}_{N}\right\}
\left\vert \Psi\right\rangle
\end{align}


\item
\begin{align}
&  \mathbf{P}\breve{\nu}_{N}\left\vert \Psi\right\rangle =\left(
-i\nabla\right)  \breve{\nu}_{N}\left\vert \Psi\right\rangle \nonumber\\
&  =\breve{\nu}_{N}\left(  -i\nabla\right)  \left\vert \Psi\right\rangle
+\left(  -i\nabla\breve{\nu}_{N}\right)  \left\vert \Psi\right\rangle
=\left\{  \breve{\nu}_{N}\mathbf{P}-i\nabla\breve{\nu}_{N}\right\}  \left\vert
\Psi\right\rangle
\end{align}


\item
\begin{equation}
\breve{\nu}_{N}e^{i\breve{\omega}_{N}}\left\vert \Psi\right\rangle
=e^{i\breve{\omega}_{N}}\breve{\nu}_{N}\left\vert \Psi\right\rangle
\end{equation}

\end{enumerate}

Which can be summarized in the commutation relationship
\begin{align}
\left[  \mathbf{P,}e^{i\breve{\omega}_{N}}\right]   &  =e^{i\breve{\omega}%
_{N}}\nabla\breve{\omega}_{N}\nonumber\\
\left[  \mathbf{P,}\breve{\nu}_{N}\right]   &  =i\nabla\breve{\nu}%
_{N}\nonumber\\
\left[  \breve{\nu}_{N},e^{i\breve{\omega}_{N}}\right]   &  =0
\end{align}
that produce
\begin{align}
\mathbf{P}\breve{\nu}_{N}e^{i\breve{\omega}_{N}}\left\vert \Psi\right\rangle
&  =\left\{  \breve{\nu}_{N}\mathbf{P}-i\nabla\breve{\nu}_{N}\right\}
e^{i\breve{\omega}_{N}}\left\vert \Psi\right\rangle =\left\{  \breve{\nu}%
_{N}e^{i\breve{\omega}_{N}}\left\{  \mathbf{P}+\nabla\breve{\omega}%
_{N}\right\}  -ie^{i\breve{\omega}_{N}}\nabla\breve{\nu}_{N}\right\}
\left\vert \Psi\right\rangle \nonumber\\
&  =e^{i\breve{\omega}_{N}}\left\{  \breve{\nu}_{N}\mathbf{P+}\breve{\nu}%
_{N}\nabla\breve{\omega}_{N}-i\nabla\breve{\nu}_{N}\right\}  \left\vert
\Psi\right\rangle
\end{align}
This allows us to express the transition amplitude%
\begin{equation}
\left\langle T\left(  \breve{\nu}_{Na},\breve{\omega}_{Na}\right)
\Psi|\left(  \mathbf{P\bullet P}\right)  T\left(  \breve{\nu}_{Nb}%
,\breve{\omega}_{Nb}\right)  \Psi\right\rangle
\end{equation}
as
\begin{equation}
\left\langle e^{i\breve{\omega}_{Na}}\left\{  \breve{\nu}_{Na}\mathbf{P+}%
\breve{\nu}_{Na}\nabla\breve{\omega}_{Na}-i\nabla\breve{\nu}_{Na}\right\}
\Psi|e^{i\breve{\omega}_{Nb}}\left\{  \breve{\nu}_{Nb}\mathbf{P+}\breve{\nu
}_{Nb}\nabla\breve{\omega}_{Nb}-i\nabla\breve{\nu}_{Nb}\right\}
\Psi\right\rangle
\end{equation}
which gives

\begin{description}
\item[A]
\begin{equation}
\left\langle \breve{\nu}_{Na}e^{i\breve{\omega}_{Na}}\mathbf{P}\Psi|\breve
{\nu}_{Nb}e^{i\breve{\omega}_{Nb}}\mathbf{P}\Psi\right\rangle =-\left\langle
\Psi|\left\{  \nabla\bullet\left(  \breve{\nu}_{Na}\breve{\nu}_{Nb}e^{i\left(
\breve{\omega}_{Nb}-\breve{\omega}_{Na}\right)  }\right)  \mathbb{I}%
\bullet\nabla\right\}  \Psi\right\rangle \label{exA}%
\end{equation}
where
\begin{align}
&  \nabla\bullet\left(  \breve{\nu}_{Na}\breve{\nu}_{Nb}e^{i\left(
\breve{\omega}_{Nb}-\breve{\omega}_{Na}\right)  }\right)  \mathbb{I}%
\bullet\nabla\nonumber\\
&  =\left(
\begin{array}
[c]{ccc}%
\nabla_{1} & \cdots & \nabla_{m}%
\end{array}
\right)  \left(
\begin{array}
[c]{ccc}%
\breve{\nu}_{Na}\breve{\nu}_{Nb}e^{i\left(  \breve{\omega}_{Nb}-\breve{\omega
}_{Na}\right)  } & \cdots & 0\\
\vdots & \ddots & \vdots\\
0 & \cdots & \breve{\nu}_{Na}\breve{\nu}_{Nb}e^{i\left(  \breve{\omega}%
_{Nb}-\breve{\omega}_{Na}\right)  }%
\end{array}
\right)  \left(
\begin{array}
[c]{c}%
\nabla_{1}\\
\vdots\\
\nabla_{m}%
\end{array}
\right)
\end{align}
The transition amplitude eq. $\left(  \ref{exA}\right)  $ is given explicitly
by
\begin{align}
&  -%
%TCIMACRO{\dint }%
%BeginExpansion
{\displaystyle\int}
%EndExpansion
\left\{  \breve{\nu}_{Na}\left(  \mathsf{z}\right)  e^{-i\breve{\omega}%
_{Na}\left(  \mathsf{z}\right)  }\nabla\Psi\left(  \mathsf{z}\right)  ^{\ast
}\breve{\nu}_{Nb}\left(  \mathsf{z}\right)  e^{i\breve{\omega}_{Nb}\left(
\mathsf{z}\right)  }\nabla\Psi\left(  \mathsf{z}\right)  \right\}
d\mathsf{z}\nonumber\\
&  \left.  =\right.  -%
%TCIMACRO{\dint }%
%BeginExpansion
{\displaystyle\int}
%EndExpansion
\left\{  \nabla\Psi\left(  \mathsf{z}\right)  ^{\dagger}\bullet\nabla
\Psi\left(  \mathsf{z}\right)  e^{i\left(  \breve{\omega}_{Nb}\left(
\mathsf{z}\right)  -\breve{\omega}_{Na}\left(  \mathsf{z}\right)  \right)
}\breve{\nu}_{Na}\left(  \mathsf{z}\right)  \breve{\nu}_{Nb}\left(
\mathsf{z}\right)  \right\}  d\mathsf{z}%
\end{align}


\item[B]
\begin{align}
&  \left\langle \breve{\nu}_{Na}e^{i\breve{\omega}_{Na}}\mathbf{P}%
\Psi|\left\{  \breve{\nu}_{Nb}\nabla\breve{\omega}_{Nb}-i\nabla\breve{\nu
}_{Nb}\right\}  e^{-i\breve{\omega}_{Nb}}\Psi\right\rangle \nonumber\\
&  =-i\left\langle \Psi|\left\{  \nabla\bullet\breve{\nu}_{Na}e^{i\left(
\breve{\omega}_{Nb}-\breve{\omega}_{Na}\right)  }\mathbb{I}\bullet\left(
\breve{\nu}_{Nb}\nabla\breve{\omega}_{Nb}-i\nabla\breve{\nu}_{Nb}\right)
\right\}  \Psi\right\rangle
\end{align}


\item[C]
\begin{align}
&  \left\langle \left(  \breve{\nu}_{Na}\nabla\breve{\omega}_{Na}%
-i\nabla\breve{\nu}_{Na}\right)  \Psi|v_{b}e^{i\breve{\omega}_{Nb}}%
\mathbf{P}\Psi\right\rangle \nonumber\\
&  =-i\left\langle \Psi|\left\{  \left(  \breve{\nu}_{Na}\nabla\breve{\omega
}_{Na}-i\nabla\breve{\nu}_{Na}\right)  \bullet v_{b}e^{i\left(  \breve{\omega
}_{Nb}-\breve{\omega}_{Na}\right)  }\mathbb{I}\bullet\nabla\right\}
\Psi\right\rangle
\end{align}


\item[D]
\begin{align}
&  \left\langle \left(  \breve{\nu}_{Na}\nabla\breve{\omega}_{Na}%
-i\nabla\breve{\nu}_{Na}\right)  \Psi|\left(  \breve{\nu}_{Nb}\nabla
\breve{\omega}_{Nb}-i\nabla\breve{\nu}_{Nb}\right)  \Psi\right\rangle
\nonumber\\
&  \left.  =\right.  \left\langle \Psi|\left\{  \left(  \breve{\nu}_{Na}%
\nabla\breve{\omega}_{Na}-i\nabla\breve{\nu}_{Na}\right)  \bullet e^{i\left(
\breve{\omega}_{Nb}-\breve{\omega}_{Na}\right)  }\mathbb{I}\bullet\left(
\breve{\nu}_{Nb}\nabla\breve{\omega}_{Nb}-i\nabla\breve{\nu}_{Nb}\right)
\right\}  \Psi\right\rangle
\end{align}
collecting together the terms A,B,C and D we obtain%
\begin{align}
&  \left\langle T\left(  \breve{\nu}_{Na},\breve{\omega}_{Na}\right)
\Psi|\left(  \mathbf{P\bullet P}\right)  T\left(  \breve{\nu}_{Nb}%
,\breve{\omega}_{Nb}\right)  \Psi\right\rangle =\nonumber\\
&  \qquad\qquad\qquad\left\langle \Psi\right.  |\left\{  -\nabla\bullet\left(
\breve{\nu}_{Na}\breve{\nu}_{Nb}e^{i\left(  \breve{\omega}_{Nb}-\breve{\omega
}_{Na}\right)  }\right)  \mathbb{I}\bullet\nabla\right. \nonumber\\
&  \qquad\qquad\qquad\qquad-\nabla\bullet\breve{\nu}_{Na}e^{i\left(
\breve{\omega}_{Nb}-\breve{\omega}_{Na}\right)  }\mathbb{I}\bullet\left(
\breve{\nu}_{Nb}\nabla\breve{\omega}_{Nb}-i\nabla\breve{\nu}_{Nb}\right)
\nonumber\\
&  \qquad\qquad\qquad\qquad\qquad-\left(  \breve{\nu}_{Na}\nabla\breve{\omega
}_{Na}-i\nabla\breve{\nu}_{Na}\right)  \bullet v_{b}e^{i\left(  \breve{\omega
}_{Nb}-\breve{\omega}_{Na}\right)  }\mathbb{I}\bullet\nabla\nonumber\\
&  \quad\left.  +\left(  \breve{\nu}_{Na}\nabla\breve{\omega}_{Na}%
-i\nabla\breve{\nu}_{Na}\right)  \bullet e^{i\left(  \breve{\omega}%
_{Nb}-\breve{\omega}_{Na}\right)  }\mathbb{I}\bullet\left(  \breve{\nu}%
_{Nb}\nabla\breve{\omega}_{Nb}-i\nabla\breve{\nu}_{Nb}\right)  \right\}
\left.  \Psi\right\rangle
\end{align}
We can also consider this as the expectation value wrt the pure state
$\left\vert \Psi\right\rangle $ of the effective operator $K\left(  \breve
{\nu}_{Na},\breve{\omega}_{Na},\breve{\nu}_{Nb},\breve{\omega}_{Nb}\right)  $
where
\begin{align}
&  K\left(  \breve{\nu}_{Na},\breve{\omega}_{Na},\breve{\nu}_{Nb}%
,\breve{\omega}_{Nb}\right)  =\nonumber\\
&  -\nabla\bullet\left(  \breve{\nu}_{Na}\breve{\nu}_{Nb}e^{i\left(
\breve{\omega}_{Nb}-\breve{\omega}_{Na}\right)  }\right)  \mathbb{I}%
\bullet\nabla-\nabla\bullet\breve{\nu}_{Na}e^{i\left(  \breve{\omega}%
_{Nb}-\breve{\omega}_{Na}\right)  }\mathbb{I}\bullet\left(  \breve{\nu}%
_{Nb}\nabla\breve{\omega}_{Nb}-i\nabla\breve{\nu}_{Nb}\right) \nonumber\\
&  \qquad\qquad\qquad-\left(  \breve{\nu}_{Na}\nabla\breve{\omega}%
_{Na}-i\nabla\breve{\nu}_{Na}\right)  \bullet v_{b}e^{i\left(  \breve{\omega
}_{Nb}-\breve{\omega}_{Na}\right)  }\mathbb{I}\bullet\nabla\nonumber\\
&  \qquad\qquad\qquad\qquad+\left(  \breve{\nu}_{Na}\nabla\breve{\omega}%
_{Na}-i\nabla\breve{\nu}_{Na}\right)  \bullet e^{i\left(  \breve{\omega}%
_{Nb}-\breve{\omega}_{Na}\right)  }\mathbb{I}\bullet\left(  \breve{\nu}%
_{Nb}\nabla\breve{\omega}_{Nb}-i\nabla\breve{\nu}_{Nb}\right)
\end{align}

\end{description}

\subsection{Expectation Values}

If we set $\breve{\nu}_{Na}=\breve{\nu}_{Nb}=\breve{\nu}_{N}$ and
$\breve{\omega}_{Na}=\breve{\omega}_{Nb}=\breve{\omega}_{N}$ we obtain the
expectation value of the Kinetic Energy operator
\begin{equation}
\left\langle T\left(  \breve{\nu}_{N},\breve{\omega}_{N}\right)  \Psi|\left(
\mathbf{P\bullet P}\right)  T\left(  \breve{\nu}_{N},\breve{\omega}%
_{N}\right)  \Psi\right\rangle
\end{equation}
which can be developed to produce
\begin{align}
&  \left\langle T\left(  \breve{\nu}_{N},\breve{\omega}_{N}\right)
\Psi|\left(  \mathbf{P\bullet P}\right)  T\left(  \breve{\nu}_{N}%
,\breve{\omega}_{N}\right)  \Psi\right\rangle =\nonumber\\
&  \qquad\qquad\qquad\qquad\left\langle \Psi\right.  |\left\{  -\nabla
\bullet\left(  \breve{\nu}_{N}^{2}\right)  \mathbb{I}\bullet\nabla\right.
\nonumber\\
&  \qquad\qquad\qquad\qquad-\nabla\bullet\left(  \breve{\nu}_{N}^{2}%
\nabla\breve{\omega}_{N}-i\breve{\nu}_{N}\nabla\breve{\nu}_{N}\right)
\nonumber\\
&  \qquad\qquad\qquad\qquad\qquad-\left(  \breve{\nu}_{N}^{2}\nabla
\breve{\omega}_{N}-i\breve{\nu}_{N}\nabla\breve{\nu}_{N}\right)  \bullet
\nabla\nonumber\\
&  \quad\left.  +\left(  \breve{\nu}_{N}\nabla\breve{\omega}_{N}-i\nabla
\breve{\nu}_{N}\right)  \bullet\mathbb{I}\bullet\left(  \breve{\nu}_{N}%
\nabla\breve{\omega}_{N}-i\nabla\breve{\nu}_{N}\right)  \right\}  \left.
\Psi\right\rangle
\end{align}
leading to the definition of an effective Kinetic Energy operator
\begin{align}
&  K\left(  \breve{\nu}_{N},\breve{\omega}_{N}\right)  =-\nabla\bullet\left(
\breve{\nu}_{N}^{2}\right)  \mathbb{I}\bullet\nabla-\nabla\bullet\left(
\breve{\nu}_{N}^{2}\nabla\breve{\omega}_{N}-i\breve{\nu}_{N}\nabla\breve{\nu
}_{N}\right) \nonumber\\
&  -\left(  \breve{\nu}_{N}^{2}\nabla\breve{\omega}_{N}-i\breve{\nu}_{N}%
\nabla\breve{\nu}_{N}\right)  \bullet\nabla+\left(  \breve{\nu}_{N}%
\nabla\breve{\omega}_{N}-i\nabla\breve{\nu}_{N}\right)  \bullet\mathbb{I}%
\bullet\left(  \breve{\nu}_{N}\nabla\breve{\omega}_{N}-i\nabla\breve{\nu}%
_{N}\right)
\end{align}


\section{Symmetric Kernels\label{symker}}

The integral kernels of the operators $\breve{T}\left(  \nu,\omega\right)  $
and in general $\left\{  \breve{T}_{n_{p}}\left(  \nu_{n_{p}},\omega_{n_{p}%
}\right)  \right\}  $ are symmetric functions. Which is shown by considering
the wavefunction
\begin{equation}
\left\langle \mathsf{z}^{n_{p}}\left\vert \tilde{\Psi}_{n_{p}}\right.
\right\rangle =\tilde{\Psi}_{n_{p}}\left(  \mathsf{z}^{n_{p}}\right)
=\nu_{n_{p}}\left(  \mathsf{z}^{n_{p}}\right)  e^{i\omega\left(
\mathsf{z}^{n_{p}}\right)  }\Psi_{n_{p}}\left(  \mathsf{z}^{n_{p}}\right)
\end{equation}
resulting from the transformation
\begin{equation}
\left\vert \tilde{\Psi}_{n_{p}}\right\rangle =\breve{T}_{n_{p}}\left(
\nu_{n_{p}},\omega_{n_{p}}\right)  \left\vert \Psi_{n_{p}}\right\rangle =\int
T_{n_{p}}\left(  \mathsf{z}^{n_{p}}\right)  \left\vert \mathsf{z}^{n_{p}%
}\right\rangle \left\langle \mathsf{z}^{n_{p}}|\Psi\right\rangle
d\mathsf{z}^{n_{p}}.
\end{equation}
$\tilde{\Psi}_{n_{p}}\left(  \mathsf{z}^{n_{p}}\right)  $ is antisymmetric wrt
coordinate permutation $\sigma$ i.e.
\begin{equation}
\tilde{\Psi}_{n_{p}}\left(  \mathsf{z}_{\sigma\left(  1\right)  }%
,\cdots,\mathsf{z}_{\sigma\left(  n_{p}\right)  }\right)  =\left(  -1\right)
^{\pi\left(  \sigma\right)  }\tilde{\Psi}_{n_{p}}\left(  \mathsf{z}^{n_{p}%
}\right)  =\left(  -1\right)  ^{\pi\left(  \sigma\right)  }\tilde{\Psi}%
_{n_{p}}\left(  \mathsf{z}_{1},\cdots,\mathsf{z}_{n_{p}}\right)  \text{ for
}\sigma\in S_{n_{p}}%
\end{equation}
where $S_{n_{p}}$ is the $\left(  n_{p}\right)  ^{th}$ order symmetric group
and $\pi\left(  \sigma\right)  $ is the parity of $\sigma$. Hence%
\begin{align}
&  \tilde{\Psi}_{n_{p}}\left(  \mathsf{z}_{\sigma\left(  1\right)  }%
,\cdots,\mathsf{z}_{\sigma\left(  n_{p}\right)  }\right)  =\left(  -1\right)
^{\pi\left(  \sigma\right)  }\tilde{\Psi}_{n_{p}}\left(  \mathsf{z}_{1}%
,\cdots,\mathsf{z}_{n_{p}}\right) \nonumber\\
&  =T_{n_{p}}\left(  \mathsf{z}_{\sigma\left(  1\right)  },\cdots
,\mathsf{z}_{\sigma\left(  n_{p}\right)  }\right)  \left(  -1\right)
^{\pi\left(  \sigma\right)  }\Psi_{n_{p}}\left(  \mathsf{z}_{1},\cdots
,\mathsf{z}_{n_{p}}\right) \nonumber\\
&  \Rightarrow T_{n_{p}}\left(  \mathsf{z}_{\sigma\left(  1\right)  }%
,\cdots,\mathsf{z}_{\sigma\left(  n_{p}\right)  }\right)  =T_{n_{p}}\left(
\mathsf{z}_{1},\cdots,\mathsf{z}_{n_{p}}\right)  \forall\sigma\in S_{n_{p}}%
\end{align}
Hence a \emph{sufficient} requirement on the magnitude $\nu_{n_{p}}\left(
\mathsf{z}_{q_{p}+1},\cdots,\mathsf{z}_{q_{p}+n_{p}}\right)  $ and phase
$\omega_{n_{p}}\left(  \mathsf{z}_{q_{p}+1},\cdots,\mathsf{z}_{q_{p}+n_{p}%
}\right)  $ to produce a symmetric $T_{n_{p}}\left(  \mathsf{z}_{q_{p}%
+1},\cdots,\mathsf{z}_{q_{p}+n_{p}}\right)  $ is that they are also symmetric
and all symmetric $T_{n_{p}}\left(  \mathsf{z}_{q_{p}+1},\cdots,\mathsf{z}%
_{q_{p}+n_{p}}\right)  $ can be generated from symmetric $\nu_{n_{p}}\left(
\mathsf{z}_{q_{p}+1},\cdots,\mathsf{z}_{q_{p}+n_{p}}\right)  $ and
$\omega_{n_{p}}\left(  \mathsf{z}_{q_{p}+1},\cdots,\mathsf{z}_{q_{p}+n_{p}%
}\right)  $, but we note that this is \emph{not a} necessary requirement.

\section{Transition Amplitudes and Expectation Values wrt IPS\label{Transamp}}

We will need to calculate expectation values and transition amplitudes between
IPS\ composed of non orthogonal $1$-particle states, to do this we will follow
the presentation \cite{Verbeek1990}. Consider two IPSs $\left\vert
\Psi\right\rangle $ and $\left\vert \Phi\right\rangle $ composed of
nonorthogonal general spin orbitals $\left\{  \psi_{j}|1\leq j\leq N\right\}
$ and $\left\{  \varphi_{j}|1\leq j\leq N\right\}  $ and the transition
amplitude
\begin{equation}
\left\langle \Psi|H\Phi\right\rangle
\end{equation}
where the Hamiltonian is a sum of one and two particle operators
\begin{equation}
H=%
%TCIMACRO{\dsum \limits_{1\leq j\leq N}}%
%BeginExpansion
{\displaystyle\sum\limits_{1\leq j\leq N}}
%EndExpansion
h\left(  j\right)  +\frac{1}{2}%
%TCIMACRO{\dsum \limits_{1\leq j\neq k\leq N}}%
%BeginExpansion
{\displaystyle\sum\limits_{1\leq j\neq k\leq N}}
%EndExpansion
g\left(  j,k\right)
\end{equation}


The transition amplitude of the $1$-particle part is given by
\begin{equation}
\left\langle \Psi\left\vert
%TCIMACRO{\dsum \limits_{1\leq j\leq N}}%
%BeginExpansion
{\displaystyle\sum\limits_{1\leq j\leq N}}
%EndExpansion
h\left(  j\right)  \right.  \Phi\right\rangle =Tr\left\{  \mathbf{h}%
\operatorname{adj}\left(  \mathbf{S}\right)  \right\}
\end{equation}
where the matrix $\mathbf{h}$ has matrix elements $\left\{  h_{jk}|1\leq
j,k\leq N\right\}  $ given by
\begin{equation}
h_{jk}=\left\langle \psi_{j}|h\left(  1\right)  \varphi_{k}\right\rangle =%
%TCIMACRO{\dint }%
%BeginExpansion
{\displaystyle\int}
%EndExpansion
\psi_{j}\left(  \mathsf{z}\right)  ^{\ast}h\left(  \mathsf{z,z}^{\prime
}\right)  \varphi_{k}\left(  \mathsf{z}^{\prime}\right)  d\mathsf{z}%
d\mathsf{z}^{\prime}%
\end{equation}
and the classical $N\times N$ adjugate matrix $\operatorname{adj}\left(
\mathbf{S}\right)  $ of the overlap matrix $\mathbf{S\equiv}\left\{
\left\langle \psi_{j}|\varphi_{k}\right\rangle |1\leq j,k\leq N\right\}  $ is
defined by
\begin{equation}
\left(  \operatorname{adj}\left(  \mathbf{S}\right)  \right)  _{jk}=\left(
-1\right)  ^{j+k}\det\left(  \mathbf{S}\left[  k|j\right]  \right)
\end{equation}
where $\mathbf{S}\left[  k|j\right]  $ is the $\left(  N-1\right)  \times$
$\left(  N-1\right)  $ matrix obtained from $\mathbf{S}$ by deleting the row
$k$ and the column $j$.

The transition amplitude of the $2$-particle part is given by
\begin{equation}
\left\langle \Psi\left\vert \frac{1}{2}%
%TCIMACRO{\dsum \limits_{1\leq j\neq k\leq N}}%
%BeginExpansion
{\displaystyle\sum\limits_{1\leq j\neq k\leq N}}
%EndExpansion
g\left(  j,k\right)  \right.  \Phi\right\rangle =\frac{1}{2}Tr\left\{
\mathbf{G}\operatorname{adj}^{\left(  2\right)  }\left(  \mathbf{S}\right)
\right\}
\end{equation}
where the matrix $\mathbf{G}$ has matrix elements $\left\{  g_{jklm}|1\leq
j<k\leq N;1\leq l<m\leq N\right\}  $ given by
\begin{align}
g_{jklm}  &  =\left\langle \psi_{j}\psi_{k}|g\left(  1,2\right)  \varphi
_{l}\varphi_{m}\right\rangle \nonumber\\
&  =%
%TCIMACRO{\dint }%
%BeginExpansion
{\displaystyle\int}
%EndExpansion
\left\{  \psi_{j}\left(  \mathsf{z}_{1}\right)  ^{\ast}\psi_{k}\left(
\mathsf{z}_{2}\right)  ^{\ast}g\left(  \mathsf{z}_{1}\mathsf{,\mathsf{z}%
_{2},\mathsf{z}_{1}^{\prime},z}_{2}^{\prime}\right)  \varphi_{l}\left(
\mathsf{z}_{1}^{\prime}\right)  \varphi_{m}\left(  \mathsf{z}_{2}^{\prime
}\right)  \right. \nonumber\\
&  \qquad\qquad\left.  -\psi_{j}\left(  \mathsf{z}_{1}\right)  ^{\ast}\psi
_{k}\left(  \mathsf{z}_{2}\right)  ^{\ast}g\left(  \mathsf{z}_{1}%
\mathsf{,\mathsf{z}_{2},\mathsf{z}_{1}^{\prime},z}_{2}^{\prime}\right)
\varphi_{m}\left(  \mathsf{z}_{1}^{\prime}\right)  \varphi_{l}\left(
\mathsf{z}_{2}^{\prime}\right)  \right\}  d\mathsf{z}^{2}d\mathsf{z}^{\prime2}%
\end{align}
and the second order $\binom{N}{2}\times\binom{N}{2}$ adjugate matrix
$\operatorname{adj}^{\left(  2\right)  }\left(  \mathbf{S}\right)  $ is
defined by
\begin{equation}
\left(  \operatorname{adj}^{\left(  2\right)  }\left(  \mathbf{S}\right)
\right)  _{jklm}=\left(  -1\right)  ^{j+k}\det\left(  \mathbf{S}\left[
lm|jk\right]  \right)
\end{equation}
where $\mathbf{S}\left[  lm|jk\right]  $ is the $\left(  N-2\right)
\times\left(  N-2\right)  $ matrix obtained from $\mathbf{S}$ by deleting the
rows $l,m$ and the columns $j,k$.

In general the transition amplitude of the $p$-particle part is given by%
\begin{equation}
\left\langle \Psi\left\vert \frac{1}{p!}%
%TCIMACRO{\dsum \limits_{1\leq j_{1}\neq\cdots\neq j_{p}\leq N}}%
%BeginExpansion
{\displaystyle\sum\limits_{1\leq j_{1}\neq\cdots\neq j_{p}\leq N}}
%EndExpansion
O_{p}\left(  j_{1},\cdots,j_{p}\right)  \right.  \Phi\right\rangle =\frac
{1}{p!}Tr\left\{  \mathbf{O}_{p}\operatorname{adj}^{\left(  p\right)  }\left(
\mathbf{S}\right)  \right\}
\end{equation}
where the matrix $\mathbf{O}_{p}$ has matrix elements
\begin{equation}
\left\{  \left[  \mathbf{O}_{p}\right]  _{j_{1},\cdots,j_{p};k_{1}%
,\cdots,k_{p}}|1\leq j_{1}<\cdots<j_{p}\leq N;1\leq k_{1}<\cdots<k_{p}\leq
N\right\}
\end{equation}
given by
\begin{align}
&  \left[  \mathbf{O}_{p}\right]  _{j_{1},\cdots,j_{p};k_{1},\cdots,k_{p}%
}=\left\langle \psi_{j_{1}}\cdots\psi_{j_{p}}|g\left(  j,k\right)
\varphi_{k_{1}}\cdots\varphi_{k_{p}}\right\rangle \nonumber\\
&  =%
%TCIMACRO{\dint }%
%BeginExpansion
{\displaystyle\int}
%EndExpansion
\left\{  \psi_{j_{1}}\left(  \mathsf{z}_{1}\right)  ^{\ast}\cdots\psi_{j_{p}%
}\left(  \mathsf{z}_{p}\right)  ^{\ast}O_{p}\left(  \mathsf{z}_{1\cdots
}\mathsf{z}_{p}\mathsf{,z}_{1\cdots}^{\prime}\mathsf{z}_{p}^{\prime}\right)
%TCIMACRO{\dsum \limits_{\sigma\in S_{p}}}%
%BeginExpansion
{\displaystyle\sum\limits_{\sigma\in S_{p}}}
%EndExpansion
\left(  -1\right)  ^{\pi\left(  \sigma\right)  }\varphi_{k_{\sigma_{\left(
1\right)  }}}\left(  \mathsf{z}_{1}^{\prime}\right)  \cdots\varphi
_{k_{\sigma\left(  p\right)  }}\left(  \mathsf{z}_{p}^{\prime}\right)
\right\}  d\mathsf{z}^{p}d\mathsf{z}^{\prime p}%
\end{align}
and the $p^{th}$ order $\binom{N}{p}\times$ $\binom{N}{p}$ adjugate matrix
$\operatorname{adj}^{\left(  p\right)  }\left(  \mathbf{S}\right)  $ is
defined by
\begin{equation}
\left(  \operatorname{adj}^{\left(  p\right)  }\left(  \mathbf{S}\right)
\right)  _{j_{1},\cdots,j_{p};k_{1},\cdots,k_{p}}=\left(  -1\right)  \,^{%
%TCIMACRO{\tsum \limits_{1\leq l\leq p}}%
%BeginExpansion
{\textstyle\sum\limits_{1\leq l\leq p}}
%EndExpansion
\left\{  j_{l}+k_{l}\right\}  }\det\left(  \mathbf{S}\left[  j_{1}%
,\cdots,j_{p}|k_{1},\cdots,k_{p}\right]  \right)
\end{equation}
where $\mathbf{S}\left[  j_{1},\cdots,j_{p}|k_{1},\cdots,k_{p}\right]  $ is
the $\left(  N-p\right)  \times\left(  N-p\right)  $ matrix obtained from
$\mathbf{S}$ by deleting the rows $j_{1},\cdots,j_{p}$ and the columns
$k_{1},\cdots,k_{p}$.

\section{ Kinetic Energy Expansion in the $\left\{  2,\ldots,2\right\}  $
Approximation\label{KinEn}}

The effective kinetic energy operator $K\left(  \nu,\omega\right)  $ is given
by%
\begin{align}
K\left(  \nu,\omega\right)   &  =\,\underset{I}{\underbrace{-\nabla
\bullet\left(  \nu^{2}\right)  \mathbb{I}\bullet\nabla}}-\underset
{II}{\underbrace{\nabla\bullet\left(  \nu^{2}\nabla\omega-i\nu\nabla
\nu\right)  }}\nonumber\\
&  -\underset{III}{\underbrace{\left(  \nu^{2}\nabla\omega-i\nu\nabla
\nu\right)  \bullet\nabla}}+\underset{IV}{\underbrace{\left(  \nu\nabla
\omega-i\nabla\nu\right)  \bullet\mathbb{I}\bullet\left(  \nu\nabla
\omega-i\nabla\nu\right)  }}%
\end{align}


\begin{description}
\item[I]
\begin{align}
&  -\nabla\bullet\left(  \nu^{2}\right)  \mathbb{I}\bullet\nabla\left.
=\right. \nonumber\\
&  -%
%TCIMACRO{\dsum \limits_{\sigma,\sigma^{\prime}\in\mathcal{S}_{N}}}%
%BeginExpansion
{\displaystyle\sum\limits_{\sigma,\sigma^{\prime}\in\mathcal{S}_{N}}}
%EndExpansion%
%TCIMACRO{\dsum \limits_{1\leq j\leq N}}%
%BeginExpansion
{\displaystyle\sum\limits_{1\leq j\leq N}}
%EndExpansion
\nabla_{j}\bullet\int\left\{
%TCIMACRO{\dprod \limits_{1\leq k\leq\frac{N}{2}}}%
%BeginExpansion
{\displaystyle\prod\limits_{1\leq k\leq\frac{N}{2}}}
%EndExpansion
\nu_{k}\left(  \mathsf{z}_{\sigma\left(  2k-1\right)  },\mathsf{z}%
_{\sigma\left(  2k\right)  }\right)  \nu_{k}\left(  \mathsf{z}_{\sigma
^{\prime}\left(  2k-1\right)  },\mathsf{z}_{\sigma^{\prime}\left(  2k\right)
}\right)  \right. \nonumber\\
&  \left.  \exp\left\{  i%
%TCIMACRO{\dsum \limits_{1\leq k\leq\frac{N}{2}}}%
%BeginExpansion
{\displaystyle\sum\limits_{1\leq k\leq\frac{N}{2}}}
%EndExpansion
\left\{  \omega_{k}\left(  \mathsf{z}_{\sigma^{\prime}\left(  2k-1\right)
},\mathsf{z}_{\sigma^{\prime}\left(  2k\right)  }\right)  -\omega_{k}\left(
\mathsf{z}_{\sigma\left(  2k-1\right)  },\mathsf{z}_{\sigma\left(  2k\right)
}\right)  \right\}  \right\}  \right\}  \left\vert \mathsf{z}^{N}\right\rangle
\left\langle \mathsf{z}^{N}\right\vert d\mathsf{z}^{N}\mathbb{I}\bullet
\nabla_{j}%
\end{align}


\item[II]
\begin{align}
&  -\nabla\bullet\left(  \nu^{2}\nabla\omega-i\nu\nabla\nu\right)  \left.
=\right. \nonumber\\
&  -%
%TCIMACRO{\dsum \limits_{\sigma,\sigma^{\prime}\in\mathcal{S}_{N}}}%
%BeginExpansion
{\displaystyle\sum\limits_{\sigma,\sigma^{\prime}\in\mathcal{S}_{N}}}
%EndExpansion%
%TCIMACRO{\dsum \limits_{1\leq j\leq N}}%
%BeginExpansion
{\displaystyle\sum\limits_{1\leq j\leq N}}
%EndExpansion
\nabla_{j}\bullet\int\left\{  \nabla_{j}\left\{
%TCIMACRO{\dsum \limits_{1\leq k\leq\frac{N}{2}}}%
%BeginExpansion
{\displaystyle\sum\limits_{1\leq k\leq\frac{N}{2}}}
%EndExpansion
\left\{  \omega_{k}\left(  \mathsf{z}_{\sigma^{\prime}\left(  2k-1\right)
},\mathsf{z}_{\sigma^{\prime}\left(  2k\right)  }\right)  \right\}  \right\}
\right.  \times\nonumber\\
&  \qquad\qquad\qquad\qquad%
%TCIMACRO{\dprod \limits_{1\leq k\leq\frac{N}{2}}}%
%BeginExpansion
{\displaystyle\prod\limits_{1\leq k\leq\frac{N}{2}}}
%EndExpansion
\nu_{k}\left(  \mathsf{z}_{\sigma\left(  2k-1\right)  },\mathsf{z}%
_{\sigma\left(  2k\right)  }\right)  \nu_{k}\left(  \mathsf{z}_{\sigma
^{\prime}\left(  2k-1\right)  },\mathsf{z}_{\sigma^{\prime}\left(  2k\right)
}\right) \nonumber\\
&  \left.  -i\nabla_{j}\left\{
%TCIMACRO{\dprod \limits_{1\leq k\leq\frac{N}{2}}}%
%BeginExpansion
{\displaystyle\prod\limits_{1\leq k\leq\frac{N}{2}}}
%EndExpansion
\nu_{k}\left(  \mathsf{z}_{\sigma^{\prime}\left(  2k-1\right)  }%
,\mathsf{z}_{\sigma^{\prime}\left(  2k\right)  }\right)  \right\}
%TCIMACRO{\dprod \limits_{1\leq k\leq\frac{N}{2}}}%
%BeginExpansion
{\displaystyle\prod\limits_{1\leq k\leq\frac{N}{2}}}
%EndExpansion
\nu_{k}\left(  \mathsf{z}_{\sigma\left(  2k-1\right)  },\mathsf{z}%
_{\sigma\left(  2k\right)  }\right)  \right\}  \left\vert \mathsf{z}%
^{N}\right\rangle \left\langle \mathsf{z}^{N}\right\vert d\mathsf{z}^{N}%
\end{align}


\item[III]
\begin{align}
&  -\left(  \nu^{2}\nabla\omega-i\nu\nabla\nu\right)  \bullet\nabla\left.
=\right. \nonumber\\
&  -%
%TCIMACRO{\dsum \limits_{\sigma,\sigma^{\prime}\in\mathcal{S}_{N}}}%
%BeginExpansion
{\displaystyle\sum\limits_{\sigma,\sigma^{\prime}\in\mathcal{S}_{N}}}
%EndExpansion%
%TCIMACRO{\dsum \limits_{1\leq j\leq N}}%
%BeginExpansion
{\displaystyle\sum\limits_{1\leq j\leq N}}
%EndExpansion
\int\left\{
%TCIMACRO{\dprod \limits_{1\leq k\leq\frac{N}{2}}}%
%BeginExpansion
{\displaystyle\prod\limits_{1\leq k\leq\frac{N}{2}}}
%EndExpansion
\nu_{k}\left(  \mathsf{z}_{\sigma\left(  2k-1\right)  },\mathsf{z}%
_{\sigma\left(  2k\right)  }\right)  \nu_{j}\left(  \mathsf{z}_{\sigma
^{\prime}\left(  2k-1\right)  },\mathsf{z}_{\sigma^{\prime}\left(  2k\right)
}\right)  \right.  \times\nonumber\\
&  \qquad\qquad\qquad\qquad\nabla_{j}\left\{
%TCIMACRO{\dsum \limits_{1\leq k\leq\frac{N}{2}}}%
%BeginExpansion
{\displaystyle\sum\limits_{1\leq k\leq\frac{N}{2}}}
%EndExpansion
\left\{  \omega_{k}\left(  \mathsf{z}_{\sigma^{\prime}\left(  2k-1\right)
},\mathsf{z}_{\sigma^{\prime}\left(  2k\right)  }\right)  \right\}  \right\}
\nonumber\\
&  \left.  -i%
%TCIMACRO{\dprod \limits_{1\leq k\leq\frac{N}{2}}}%
%BeginExpansion
{\displaystyle\prod\limits_{1\leq k\leq\frac{N}{2}}}
%EndExpansion
\nu_{k}\left(  \mathsf{z}_{\sigma\left(  2k-1\right)  },\mathsf{z}%
_{\sigma\left(  2k\right)  }\right)  \nabla_{j}\left\{
%TCIMACRO{\dprod \limits_{1\leq k\leq\frac{N}{2}}}%
%BeginExpansion
{\displaystyle\prod\limits_{1\leq k\leq\frac{N}{2}}}
%EndExpansion
\nu_{k}\left(  \mathsf{z}_{\sigma^{\prime}\left(  2k-1\right)  }%
,\mathsf{z}_{\sigma^{\prime}\left(  2k\right)  }\right)  \right\}  \right\}
\left\vert \mathsf{z}^{N}\right\rangle \left\langle \mathsf{z}^{\prime
N}\right\vert d\mathsf{z}^{N}\bullet\nabla_{j}%
\end{align}


\item[IV]
\begin{align}
&  \left(  \nu\nabla\omega-i\nabla\nu\right)  \bullet\mathbb{I}\bullet\left(
\nu\nabla\omega-i\nabla\nu\right)  \left.  =\right. \nonumber\\
&
%TCIMACRO{\dsum \limits_{\sigma,\sigma^{\prime}\in\mathcal{S}_{N}}}%
%BeginExpansion
{\displaystyle\sum\limits_{\sigma,\sigma^{\prime}\in\mathcal{S}_{N}}}
%EndExpansion%
%TCIMACRO{\dsum \limits_{1\leq j\leq N}}%
%BeginExpansion
{\displaystyle\sum\limits_{1\leq j\leq N}}
%EndExpansion
\int%
%TCIMACRO{\dprod \limits_{1\leq k\leq\frac{N}{2}}}%
%BeginExpansion
{\displaystyle\prod\limits_{1\leq k\leq\frac{N}{2}}}
%EndExpansion
\nu_{k}\left(  \mathsf{z}_{\sigma\left(  2k-1\right)  },\mathsf{z}%
_{\sigma\left(  2k\right)  }\right)  \times\nonumber\\
&  \left\{  \nabla_{j}\left\{
%TCIMACRO{\dsum \limits_{1\leq k\leq\frac{N}{2}}}%
%BeginExpansion
{\displaystyle\sum\limits_{1\leq k\leq\frac{N}{2}}}
%EndExpansion
\left\{  \omega_{k}\left(  \mathsf{z}_{\sigma\left(  2k-1\right)  }%
,\mathsf{z}_{\sigma\left(  2k\right)  }\right)  \right\}  \right\}
-i\nabla_{j}\left\{
%TCIMACRO{\dprod \limits_{1\leq k\leq\frac{N}{2}}}%
%BeginExpansion
{\displaystyle\prod\limits_{1\leq k\leq\frac{N}{2}}}
%EndExpansion
\nu_{k}\left(  \mathsf{z}_{\sigma\left(  2k-1\right)  },\mathsf{z}%
_{\sigma\left(  2k\right)  }\right)  \right\}  \right\}  \left\vert
\mathsf{z}^{N}\right\rangle \left\langle \mathsf{z}^{N}\right\vert
d\mathsf{z}^{N}\nonumber\\
&  \bullet\mathbb{I}\bullet%
%TCIMACRO{\dsum \limits_{\sigma,\sigma^{\prime}\in\mathcal{S}_{N}}}%
%BeginExpansion
{\displaystyle\sum\limits_{\sigma,\sigma^{\prime}\in\mathcal{S}_{N}}}
%EndExpansion%
%TCIMACRO{\dsum \limits_{1\leq j\leq N}}%
%BeginExpansion
{\displaystyle\sum\limits_{1\leq j\leq N}}
%EndExpansion
\int\left\{
%TCIMACRO{\dprod \limits_{1\leq k\leq\frac{N}{2}}}%
%BeginExpansion
{\displaystyle\prod\limits_{1\leq k\leq\frac{N}{2}}}
%EndExpansion
\nu_{k}\left(  \mathsf{z}_{\sigma^{\prime}\left(  2k-1\right)  }%
,\mathsf{z}_{\sigma^{\prime}\left(  2k\right)  }\right)  \right.  \times\\
&  \nabla_{j}\left\{
%TCIMACRO{\dsum \limits_{1\leq k\leq\frac{N}{2}}}%
%BeginExpansion
{\displaystyle\sum\limits_{1\leq k\leq\frac{N}{2}}}
%EndExpansion
\left\{  \omega_{k}\left(  \mathsf{z}_{\sigma^{\prime}\left(  2k-1\right)
},\mathsf{z}_{\sigma^{\prime}\left(  2k\right)  }\right)  \right\}  \right\}
\left.  -i\nabla_{j}\left\{
%TCIMACRO{\dprod \limits_{1\leq k\leq\frac{N}{2}}}%
%BeginExpansion
{\displaystyle\prod\limits_{1\leq k\leq\frac{N}{2}}}
%EndExpansion
\nu_{k}\left(  \mathsf{z}_{\sigma^{\prime}\left(  2k-1\right)  }%
,\mathsf{z}_{\sigma^{\prime}\left(  2k\right)  }\right)  \right\}  \right\}
\left\vert \mathsf{z}^{N}\right\rangle \left\langle \mathsf{z}^{N}\right\vert
d\mathsf{z}^{N}%
\end{align}

\end{description}

leading to the expectation values

\begin{description}
\item[I]
\begin{align}
&  \mathcal{K}\left(  \nu,\omega\right)  =\left\langle \varphi_{1}%
\cdots\varphi_{N}|K\left(  \nu,\omega\right)  \varphi_{1}\cdots\varphi
_{N}\right\rangle =\nonumber\\
&
%TCIMACRO{\dsum \limits_{\sigma,\sigma^{\prime},\sigma^{\prime\prime}%
%\in\mathcal{S}_{N}}}%
%BeginExpansion
{\displaystyle\sum\limits_{\sigma,\sigma^{\prime},\sigma^{\prime\prime}%
\in\mathcal{S}_{N}}}
%EndExpansion%
%TCIMACRO{\dint }%
%BeginExpansion
{\displaystyle\int}
%EndExpansion
\left\{  \left\{  \varphi_{_{1}}\left(  \mathsf{z}_{1}\right)  ^{\ast}%
\cdots\nabla_{j}\varphi_{j}\left(  \mathsf{z}_{j}\right)  ^{\ast}\cdots
\varphi_{N}\left(  \mathsf{z}_{N}\right)  ^{\ast}\right\}  \bullet\nabla
_{j}\left\{  \left(  -1\right)  ^{\pi\left(  \sigma^{\prime\prime}\right)
}\varphi_{1}\left(  \mathsf{z}_{\sigma^{\prime\prime}\left(  1\right)
}\right)  \cdots\varphi_{N}\left(  \mathsf{z}_{\sigma^{\prime\prime}\left(
N\right)  }\right)  \right\}  \right. \nonumber\\
&  \left.
%TCIMACRO{\dprod \limits_{1\leq k\leq\frac{N}{2}}}%
%BeginExpansion
{\displaystyle\prod\limits_{1\leq k\leq\frac{N}{2}}}
%EndExpansion
\nu_{k}\left(  \mathsf{z}_{\sigma\left(  2k-1\right)  },\mathsf{z}%
_{\sigma\left(  2k\right)  }\right)  \nu_{k}\left(  \mathsf{z}_{\sigma
^{\prime}\left(  2k-1\right)  },\mathsf{z}_{\sigma^{\prime}\left(  2k\right)
}\right)  e^{i%
%TCIMACRO{\tsum \limits_{1\leq k\leq\frac{N}{2}}}%
%BeginExpansion
{\textstyle\sum\limits_{1\leq k\leq\frac{N}{2}}}
%EndExpansion
\left\{  \omega_{k}\left(  \mathsf{z}_{\sigma^{\prime}\left(  2k-1\right)
},\mathsf{z}_{\sigma^{\prime}\left(  2k\right)  }\right)  -\omega_{k}\left(
\mathsf{z}_{\sigma\left(  2k-1\right)  },\mathsf{z}_{\sigma\left(  2k\right)
}\right)  \right\}  }\right\}  d\mathsf{z}^{N}%
\end{align}


\item[II]
\begin{align}
&  -%
%TCIMACRO{\dsum \limits_{\sigma,\sigma^{\prime},\sigma^{\prime\prime}%
%\in\mathcal{S}_{N}}}%
%BeginExpansion
{\displaystyle\sum\limits_{\sigma,\sigma^{\prime},\sigma^{\prime\prime}%
\in\mathcal{S}_{N}}}
%EndExpansion%
%TCIMACRO{\dsum \limits_{1\leq j\leq N}}%
%BeginExpansion
{\displaystyle\sum\limits_{1\leq j\leq N}}
%EndExpansion
\int\left\{  \left\{  \left(  -1\right)  ^{\pi\left(  \sigma^{\prime\prime
}\right)  }\varphi_{1}\left(  \mathsf{z}_{\sigma^{\prime\prime}\left(
1\right)  }\right)  \cdots\varphi_{N}\left(  \mathsf{z}_{\sigma^{\prime\prime
}\left(  N\right)  }\right)  \right\}  \right. \nonumber\\
&  \varphi_{_{1}}\left(  \mathsf{z}_{1}\right)  ^{\ast}\cdots\nabla_{j}%
\varphi_{j}\left(  \mathsf{z}_{j}\right)  ^{\ast}\cdots\varphi_{N}\left(
\mathsf{z}_{N}\right)  ^{\ast}\bullet\left\{  \nabla_{j}\left\{
%TCIMACRO{\dsum \limits_{1\leq k\leq\frac{N}{2}}}%
%BeginExpansion
{\displaystyle\sum\limits_{1\leq k\leq\frac{N}{2}}}
%EndExpansion
\left\{  \omega_{k}\left(  \mathsf{z}_{\sigma^{\prime}\left(  2k-1\right)
},\mathsf{z}_{\sigma^{\prime}\left(  2k\right)  }\right)  \right\}  \right\}
\right. \nonumber\\
&  \qquad\qquad\left.  \qquad\right.  \qquad%
%TCIMACRO{\dprod \limits_{1\leq k\leq\frac{N}{2}}}%
%BeginExpansion
{\displaystyle\prod\limits_{1\leq k\leq\frac{N}{2}}}
%EndExpansion
\nu_{k}\left(  \mathsf{z}_{\sigma\left(  2k-1\right)  },\mathsf{z}%
_{\sigma\left(  2k\right)  }\right)  \nu_{k}\left(  \mathsf{z}_{\sigma
^{\prime}\left(  2k-1\right)  },\mathsf{z}_{\sigma^{\prime}\left(  2k\right)
}\right) \nonumber\\
&  \qquad\qquad\left.  -\,i\nabla_{j}\left\{
%TCIMACRO{\dprod \limits_{1\leq k\leq\frac{N}{2}}}%
%BeginExpansion
{\displaystyle\prod\limits_{1\leq k\leq\frac{N}{2}}}
%EndExpansion
\nu_{k}\left(  \mathsf{z}_{\sigma^{\prime}\left(  2k-1\right)  }%
,\mathsf{z}_{\sigma^{\prime}\left(  2k\right)  }\right)  \right\}
%TCIMACRO{\dprod \limits_{1\leq k\leq\frac{N}{2}}}%
%BeginExpansion
{\displaystyle\prod\limits_{1\leq k\leq\frac{N}{2}}}
%EndExpansion
\nu_{k}\left(  \mathsf{z}_{\sigma\left(  2k-1\right)  },\mathsf{z}%
_{\sigma\left(  2k\right)  }\right)  \right\}  d\mathsf{z}^{N}%
\end{align}


\item[III]
\begin{align}
&  -%
%TCIMACRO{\dsum \limits_{\sigma,\sigma^{\prime},\sigma^{\prime\prime}%
%\in\mathcal{S}_{N}}}%
%BeginExpansion
{\displaystyle\sum\limits_{\sigma,\sigma^{\prime},\sigma^{\prime\prime}%
\in\mathcal{S}_{N}}}
%EndExpansion%
%TCIMACRO{\dsum \limits_{1\leq j\leq N}}%
%BeginExpansion
{\displaystyle\sum\limits_{1\leq j\leq N}}
%EndExpansion
\int\varphi_{_{1}}\left(  \mathsf{z}_{1}\right)  ^{\ast}\cdots\varphi
_{N}\left(  \mathsf{z}_{N}\right)  ^{\ast}\left\{
%TCIMACRO{\dprod \limits_{1\leq k\leq\frac{N}{2}}}%
%BeginExpansion
{\displaystyle\prod\limits_{1\leq k\leq\frac{N}{2}}}
%EndExpansion
\nu_{k}\left(  \mathsf{z}_{\sigma\left(  2k-1\right)  },\mathsf{z}%
_{\sigma\left(  2k\right)  }\right)  \nu_{k}\left(  \mathsf{z}_{\sigma
^{\prime}\left(  2k-1\right)  },\mathsf{z}_{\sigma^{\prime}\left(  2k\right)
}\right)  \right. \nonumber\\
&  \qquad\qquad\qquad\times%
%TCIMACRO{\dsum \limits_{1\leq k\leq\frac{N}{2}}}%
%BeginExpansion
{\displaystyle\sum\limits_{1\leq k\leq\frac{N}{2}}}
%EndExpansion
\nabla_{j}\left\{  \omega_{k}\left(  \mathsf{z}_{\sigma^{\prime}\left(
2k-1\right)  },\mathsf{z}_{\sigma^{\prime}\left(  2k\right)  }\right)
\right\}  -i%
%TCIMACRO{\dprod \limits_{1\leq k\leq\frac{N}{2}}}%
%BeginExpansion
{\displaystyle\prod\limits_{1\leq k\leq\frac{N}{2}}}
%EndExpansion
\nu_{k}\left(  \mathsf{z}_{\sigma\left(  2k-1\right)  },\mathsf{z}%
_{\sigma\left(  2k\right)  }\right)  \times\nonumber\\
&  \nabla_{j}\left\{
%TCIMACRO{\dprod \limits_{1\leq k\leq\frac{N}{2}}}%
%BeginExpansion
{\displaystyle\prod\limits_{1\leq k\leq\frac{N}{2}}}
%EndExpansion
\nu_{k}\left(  \mathsf{z}_{\sigma^{\prime}\left(  2k-1\right)  }%
,\mathsf{z}_{\sigma^{\prime}\left(  2k\right)  }\right)  \right\}  \left.
\bullet\nabla_{j}\left\{  \left(  -1\right)  ^{\pi\left(  \sigma^{\prime
\prime}\right)  }\varphi_{1}\left(  \mathsf{z}_{\sigma^{\prime\prime}\left(
1\right)  }\right)  \cdots\varphi_{N}\left(  \mathsf{z}_{\sigma^{\prime\prime
}\left(  N\right)  }\right)  \right\}  \right\}  d\mathsf{z}^{N}%
\end{align}


\item[IV]
\begin{align}
&  +%
%TCIMACRO{\dsum \limits_{\sigma,\sigma^{\prime},\sigma^{\prime\prime}%
%\in\mathcal{S}_{N}}}%
%BeginExpansion
{\displaystyle\sum\limits_{\sigma,\sigma^{\prime},\sigma^{\prime\prime}%
\in\mathcal{S}_{N}}}
%EndExpansion%
%TCIMACRO{\dsum \limits_{1\leq j\leq N}}%
%BeginExpansion
{\displaystyle\sum\limits_{1\leq j\leq N}}
%EndExpansion
\int\varphi_{_{1}}\left(  \mathsf{z}_{1}\right)  ^{\ast}\cdots\varphi
_{N}\left(  \mathsf{z}_{N}\right)  ^{\ast}\left(  -1\right)  ^{\pi\left(
\sigma^{\prime\prime}\right)  }\varphi_{1}\left(  \mathsf{z}_{\sigma
^{\prime\prime}\left(  1\right)  }\right)  \cdots\varphi_{N}\left(
\mathsf{z}_{\sigma^{\prime\prime}\left(  N\right)  }\right)  \times\nonumber\\
&  \left\{
%TCIMACRO{\dprod \limits_{1\leq k\leq\frac{N}{2}}}%
%BeginExpansion
{\displaystyle\prod\limits_{1\leq k\leq\frac{N}{2}}}
%EndExpansion
\nu_{k}\left(  \mathsf{z}_{\sigma\left(  2k-1\right)  },\mathsf{z}%
_{\sigma\left(  2k\right)  }\right)  \nabla_{j}\left\{
%TCIMACRO{\dsum \limits_{1\leq k\leq\frac{N}{2}}}%
%BeginExpansion
{\displaystyle\sum\limits_{1\leq k\leq\frac{N}{2}}}
%EndExpansion
\left\{  \omega_{k}\left(  \mathsf{z}_{\sigma\left(  2k-1\right)  }%
,\mathsf{z}_{\sigma\left(  2k\right)  }\right)  \right\}  \right\}  \right.
\nonumber\\
&  -i\nabla_{j}\left\{
%TCIMACRO{\dprod \limits_{1\leq k\leq\frac{N}{2}}}%
%BeginExpansion
{\displaystyle\prod\limits_{1\leq k\leq\frac{N}{2}}}
%EndExpansion
\nu_{k}\left(  \mathsf{z}_{\sigma\left(  2k-1\right)  },\mathsf{z}%
_{\sigma\left(  2k\right)  }\right)  \right\}  \bullet\left\{  \nabla
_{j}\left\{
%TCIMACRO{\dsum \limits_{1\leq k\leq\frac{N}{2}}}%
%BeginExpansion
{\displaystyle\sum\limits_{1\leq k\leq\frac{N}{2}}}
%EndExpansion
\left\{  \omega_{k}\left(  \mathsf{z}_{\sigma^{\prime}\left(  2k-1\right)
},\mathsf{z}_{\sigma^{\prime}\left(  2k\right)  }\right)  \right\}  \right\}
\right.  \times\nonumber\\
&  \qquad\qquad%
%TCIMACRO{\dprod \limits_{1\leq k\leq\frac{N}{2}}}%
%BeginExpansion
{\displaystyle\prod\limits_{1\leq k\leq\frac{N}{2}}}
%EndExpansion
\nu_{k}\left(  \mathsf{z}_{\sigma^{\prime}\left(  2k-1\right)  }%
,\mathsf{z}_{\sigma^{\prime}\left(  2k\right)  }\right)  \left.  -i\nabla
_{j}\left\{
%TCIMACRO{\dprod \limits_{1\leq k\leq\frac{N}{2}}}%
%BeginExpansion
{\displaystyle\prod\limits_{1\leq k\leq\frac{N}{2}}}
%EndExpansion
\nu_{k}\left(  \mathsf{z}_{\sigma^{\prime}\left(  2k-1\right)  }%
,\mathsf{z}_{\sigma^{\prime}\left(  2k\right)  }\right)  \right\}  \right\}
d\mathsf{z}^{N}%
\end{align}

\end{description}

\bibliographystyle{apsrev}
\bibliography{novel}

\end{document}